% Blueprint content for the Offer acceptance slice.
%
% Do not write \leanok here.  scripts/render_status.py reads build/status.json
% and writes content.generated.tex with \leanok added to every environment
% whose declaration the build reports as defined or proven.

\chapter{Offer acceptance}

This blueprint covers one slice of the Offer and Order domain: an offer with
priced items and an expiry, the order created when the offer is accepted,
payments against the order, and refunds against payments.  Each node in this
chapter is a Lean declaration.  A green node is defined or proven in Lean; an
uncolored theorem is open.  The next chapter adds candidate rules from IATA's
public ONE Order material that have no Lean yet.

The review document, under \texttt{review/} on the site, shows the same declarations with
their docstrings, the examples the build computes, and the open question for
the Offer/Order team.

\section{Core concepts}

\begin{definition}[Currency]
  \label{def:currency}
  \lean{Rules.Core.Currency}
  The currencies the slice prices offers in, as ISO 4217 codes: USD, EUR and
  GBP.
\end{definition}

\begin{definition}[Money]
  \label{def:money}
  \lean{Rules.Core.Money}
  \uses{def:currency}
  An amount of money is a non-negative whole number of minor units (cents,
  pence) in one currency.
\end{definition}

\begin{definition}[Adding money]
  \label{def:money-add}
  \lean{Rules.Core.Money.add}
  \uses{def:money}
  Two amounts add only when they are in the same currency.  A mixed-currency
  sum has no result, so it is a visible failure rather than a silent
  conversion.
\end{definition}

\begin{definition}[Instant]
  \label{def:instant}
  \lean{Rules.Core.Instant}
  A point in time, as whole seconds since the Unix epoch.  All times are UTC.
\end{definition}

\begin{definition}[Passenger]
  \label{def:passenger}
  \lean{Rules.Core.Passenger}
  A traveller named on an offer or order.
\end{definition}

\begin{definition}[Principal]
  \label{def:principal}
  \lean{Rules.Core.Principal}
  Whoever makes a request.  A placeholder for the Identity namespace.
\end{definition}

\begin{definition}[Classification]
  \label{def:classification}
  \lean{Rules.Core.Classification}
  Data sensitivity levels, ordered public, internal, confidential,
  restricted.  A placeholder for the Data Governance namespace.
\end{definition}

\begin{theorem}[Public is the lowest level]
  \label{thm:public-le}
  \lean{Rules.Core.Classification.public_le}
  \uses{def:classification}
  Every classification level is at least as sensitive as public.
\end{theorem}
\begin{proof}
  By comparing ranks: public has rank 0.
\end{proof}

\section{Offers and orders}

\begin{definition}[Offer item]
  \label{def:offer-item}
  \lean{Rules.OfferOrder.OfferItem}
  \uses{def:money, def:passenger}
  One priced product inside an offer for one passenger: a flight, a bag, a
  seat.  It corresponds to a ONE Order OfferItem.
\end{definition}

\begin{definition}[Offer]
  \label{def:offer}
  \lean{Rules.OfferOrder.Offer}
  \uses{def:offer-item, def:instant, def:currency}
  A set of priced items the airline offers a customer, valid until it
  expires, with every price in one currency.
\end{definition}

\begin{definition}[Delivery state]
  \label{def:delivery-state}
  \lean{Rules.OfferOrder.DeliveryState}
  Where a purchased service is in its life: ready, delivered, cancelled, or
  refunded.
\end{definition}

\begin{definition}[Delivery event]
  \label{def:delivery-event}
  \lean{Rules.OfferOrder.DeliveryEvent}
  Something that happens to a purchased service: deliver, cancel, or refund.
\end{definition}

\begin{definition}[Order item]
  \label{def:order-item}
  \lean{Rules.OfferOrder.OrderItem}
  \uses{def:money, def:delivery-state}
  One purchased service inside an order, created from one offer item.
\end{definition}

\begin{definition}[Order]
  \label{def:order}
  \lean{Rules.OfferOrder.Order}
  \uses{def:order-item, def:money, def:instant}
  A customer's purchase: the items of one accepted offer, the total price, and
  when it was created.
\end{definition}

\begin{definition}[Rejection]
  \label{def:rejection}
  \lean{Rules.OfferOrder.Rejection}
  Why an offer could not be accepted: it expired, it is empty, or its items
  are priced in more than one currency.
\end{definition}

\begin{definition}[Accepting an offer]
  \label{def:accept}
  \lean{Rules.OfferOrder.accept}
  \uses{def:offer, def:order, def:rejection}
  When a customer accepts an offer, the system creates an order only if the
  offer has not expired, contains at least one item, and prices every item in
  the offer's currency.  The order's price is the sum of the item prices.
\end{definition}

\begin{definition}[Delivery transitions]
  \label{def:step}
  \lean{Rules.OfferOrder.step}
  \uses{def:delivery-state, def:delivery-event}
  The only allowed transitions are ready to delivered, ready to cancelled, and
  cancelled to refunded.  Every other event is refused.
\end{definition}

\begin{theorem}[Order price is the sum of the items]
  \label{thm:accept-price-eq-sum}
  \lean{Rules.OfferOrder.accept_price_eq_sum}
  \uses{def:accept}
  When an offer is accepted, the order's price is exactly the sum of the
  offer's item prices.
\end{theorem}
\begin{proof}
  \uses{def:accept}
  Unfold the acceptance decision; in the only branch that creates an order,
  the price is defined as that sum.
\end{proof}

\begin{theorem}[Expired offers are rejected]
  \label{thm:accept-rejects-expired}
  \lean{Rules.OfferOrder.accept_rejects_expired}
  \uses{def:accept}
  An offer accepted at or after its expiry instant is rejected as expired.
\end{theorem}
\begin{proof}
  \uses{def:accept}
  The expiry check is the first check the decision makes.
\end{proof}

\begin{theorem}[Delivered services cannot be cancelled]
  \label{thm:delivered-not-cancellable}
  \lean{Rules.OfferOrder.delivered_not_cancellable}
  \uses{def:step}
  The cancel event is refused for a delivered service, and no event moves a
  delivered service to cancelled.
\end{theorem}
\begin{proof}
  \uses{def:step}
  Check each of the three events.
\end{proof}

\section{Payments and refunds}

\begin{definition}[Payment]
  \label{def:payment}
  \lean{Rules.OfferOrder.Payment}
  \uses{def:money}
  Money received from the customer against an order.
\end{definition}

\begin{definition}[Refund]
  \label{def:refund}
  \lean{Rules.OfferOrder.Refund}
  \uses{def:money}
  Money returned to the customer against one payment.
\end{definition}

\begin{definition}[Totals paid and refunded]
  \label{def:totals}
  \lean{Rules.OfferOrder.totalPaid, Rules.OfferOrder.totalRefunded}
  \uses{def:payment, def:refund}
  The total received in an order's payments, and the total returned in
  refunds against those payments.
\end{definition}

\begin{definition}[Accepting a refund]
  \label{def:accept-refund}
  \lean{Rules.OfferOrder.acceptRefund}
  \uses{def:payment, def:refund}
  A refund is accepted against a payment only if the refund amount does not
  exceed that payment's amount.  The check looks at one refund at a time.
\end{definition}

\begin{theorem}[Refunds never exceed payments]
  \label{thm:refunds-le-payments}
  \lean{Rules.OfferOrder.refunds_le_payments}
  \uses{def:accept-refund, def:totals}
  If every refund was accepted against one of the order's payments, the total
  refunded does not exceed the total paid.

  Counterexample: two refunds of 60.00 against one payment of 100.00 are each
  accepted and together return 120.00.

  Open point: does the limit apply per refund, cumulatively per payment, or
  cumulatively per order, and is a refund above the amount paid ever allowed,
  for example as a goodwill credit?
\end{theorem}
\begin{proof}
  \uses{def:accept-refund, def:totals}
  Not proven.  The counterexample above shows the statement is false for the
  rule as written.
\end{proof}

\input{ears.generated}

\chapter{ONE Order rules from the survey}

This chapter lists the candidate rules, by area of the ONE Order coverage map, chosen in the
\href{https://github.com/samdengler/airline-order-verified/blob/main/docs/one-order-survey.md}{ONE Order survey}.
None of them is in Lean yet.  Each node describes a rule in our own words and
links to the public IATA page it comes from; IATA owns the standard and its
documentation.  A node marked not ready is one where the source is vague or
silent, so the business has to decide what the rule means before it can be
written in Lean.  The open point in each such node says what needs deciding.

% Candidate rules for order structure and status.  Written by hand; see docs/one-order-coverage.md.

\section{Order structure and status}

\begin{definition}[Order structure (C-01)]
  \label{def:c01-structure}
  \lean{Rules.OneOrder.Service, Rules.OneOrder.ServiceStatus, Rules.OneOrder.OrderItem, Rules.OneOrder.Order}
  \uses{def:order, def:order-item, def:passenger}
  An order holds one or more order items, and each order item holds one or
  more services.  A service is for one passenger on one flight segment, and is
  the smallest part of an order that can be delivered on its own.  Each
  service has a service status, such as \texttt{CONFIRMED} or
  \texttt{TRANSFERRED}, and a delivery status, such as
  \texttt{READY TO PROCEED}.

  Source: \href{https://standards.atlassian.net/wiki/spaces/EASD/pages/574849062/Order}{EASD guide, Order}.
\end{definition}

\begin{definition}[Offer items for passengers and segments]
  \label{def:one-offer}
  \lean{Rules.OneOrder.OfferItem, Rules.OneOrder.Offer}
  \uses{def:offer, def:passenger, def:money, def:instant}
  An offer item is priced for one or more passengers on one or more flight
  segments.  An offer holds its items, an expiry instant and a currency.

  Source: \href{https://standards.atlassian.net/wiki/spaces/EASD/pages/574849062/Order}{EASD guide, Order}.
\end{definition}

\begin{definition}[Creating an order from an offer]
  \label{def:one-accept}
  \lean{Rules.OneOrder.accept}
  \uses{def:one-offer, def:c01-structure, def:rejection}
  Accepting an offer before it expires creates an order with one order item
  for each offer item, and one confirmed service for each passenger and
  segment of that item.  The checks before creating the order are the same as
  in the pilot rule: expiry, then an empty offer, then mixed currencies.

  Source: \href{https://standards.atlassian.net/wiki/spaces/EASD/pages/574849062/Order}{EASD guide, Order}.
\end{definition}

\begin{theorem}[One order item per offer item (C-02)]
  \label{thm:c02-item-per-offer-item}
  \lean{Rules.OneOrder.accept_item_per_offer_item}
  \uses{def:one-accept, def:c01-structure}
  Accepting an offer creates exactly one order item for each offer item.

  Source: \href{https://standards.atlassian.net/wiki/spaces/EASD/pages/574849062/Order}{EASD guide, Order}.
\end{theorem}
\begin{proof}
  \uses{def:one-accept}
  The order's items are the offer's items, each turned into an order item that
  records where it came from.
\end{proof}

\begin{theorem}[One service per passenger and segment]
  \label{thm:one-services-per-passenger-segment}
  \lean{Rules.OneOrder.accept_services_per_passenger_segment}
  \uses{def:one-accept, def:c01-structure}
  Each order item created from an offer item holds one service for each
  passenger and flight segment of that offer item.

  Source: \href{https://standards.atlassian.net/wiki/spaces/EASD/pages/574849062/Order}{EASD guide, Order}.
\end{theorem}
\begin{proof}
  \uses{def:one-accept}
  The services are built from the pairs of a passenger and a segment.
\end{proof}

\begin{theorem}[Order item IDs are unique (C-03)]
  \label{thm:c03-unique-item-ids}
  \lean{Rules.OneOrder.accept_item_ids_unique}
  \uses{def:one-accept, def:c01-structure}
  Accepting an offer whose item IDs are distinct gives the order's items
  distinct IDs.  Whether every change keeps them distinct belongs with the
  change rules (C-18).

  Source: \href{https://standards.atlassian.net/wiki/spaces/EASD/pages/574849062/Order}{EASD guide, Order}.
\end{theorem}
\begin{proof}
  \uses{def:one-accept}
  Each order item ID is the order ID, a hyphen and the offer item ID, so
  distinct offer item IDs give distinct order item IDs.
\end{proof}

\begin{theorem}[Order items are not split (C-04)]
  \label{thm:c04-item-unsplit}
  \uses{def:c01-structure, def:c18-change}
  The price is held on the order item, and no operation changes part of an
  item.  A change replaces whole items.

  Source: \href{https://standards.atlassian.net/wiki/spaces/EASD/pages/574849062/Order}{EASD guide, Order}.
\end{theorem}

\begin{definition}[Services that hold value (C-05)]
  \label{def:c05-holds-value}
  \lean{Rules.OneOrder.ValueReading}
  \uses{def:c08-delivery-status}
  Whether a service still holds value, decided by its delivery status.  The
  source settles five statuses: \texttt{CONFIRMED}, \texttt{READY TO PROCEED}
  and \texttt{REMOVED HOLD VALUE} hold value, and \texttt{DELIVERED} and
  \texttt{REMOVED} do not.  The model takes the answer for the other eight as a
  parameter, a reading of the guide, and a guarantee about value counts as
  proven only when it holds for every reading.

  Open point: which of the other eight delivery statuses hold value, for
  example \texttt{EXPIRED}, \texttt{NOT CLAIMED} and \texttt{SUSPENDED}.

  Source: \href{https://standards.atlassian.net/wiki/spaces/EASD/pages/574849062/Order}{EASD guide, Order}.
\end{definition}

\begin{definition}[Order item status (C-06)]
  \label{def:c06-item-status}
  \lean{Rules.OneOrder.itemStatuses}
  \uses{def:c01-structure, def:c05-holds-value}
  An order item is \texttt{ACTIVE} while at least one of its services holds
  value.  It is \texttt{CANCELLED} once all its services are cancelled or
  transferred and none holds value.  It is \texttt{CLOSED} once at least one
  service is delivered and every other service is delivered or removed.

  Source: \href{https://standards.atlassian.net/wiki/spaces/EASD/pages/574849062/Order}{EASD guide, Order}.
\end{definition}

\begin{theorem}[Active items are those holding value]
  \label{thm:one-item-active}
  \lean{Rules.OneOrder.item_active_iff}
  \uses{def:c06-item-status}
  An order item is marked active exactly when at least one of its services
  holds value.

  Source: \href{https://standards.atlassian.net/wiki/spaces/EASD/pages/574849062/Order}{EASD guide, Order}.
\end{theorem}
\begin{proof}
  \uses{def:c06-item-status}
  Unfold the status rule.
\end{proof}

\begin{theorem}[Cancelled items as the guide defines them]
  \label{thm:one-item-cancelled}
  \lean{Rules.OneOrder.item_cancelled_iff}
  \uses{def:c06-item-status}
  An order item is marked cancelled exactly when every service is cancelled
  or transferred and none holds value.

  Source: \href{https://standards.atlassian.net/wiki/spaces/EASD/pages/574849062/Order}{EASD guide, Order}.
\end{theorem}
\begin{proof}
  \uses{def:c06-item-status}
  Unfold the status rule.
\end{proof}

\begin{theorem}[Closed items as the guide defines them]
  \label{thm:one-item-closed}
  \lean{Rules.OneOrder.item_closed_iff}
  \uses{def:c06-item-status}
  An order item is marked closed exactly when at least one service is
  delivered and every service is delivered or removed.

  Source: \href{https://standards.atlassian.net/wiki/spaces/EASD/pages/574849062/Order}{EASD guide, Order}.
\end{theorem}
\begin{proof}
  \uses{def:c06-item-status}
  Unfold the status rule.
\end{proof}

\begin{theorem}[Active items are never closed]
  \label{thm:one-active-not-closed}
  \lean{Rules.OneOrder.active_not_closed}
  \uses{def:c06-item-status}
  Under every reading of the guide, an active item is not closed.

  Source: \href{https://standards.atlassian.net/wiki/spaces/EASD/pages/574849062/Order}{EASD guide, Order}.
\end{theorem}
\begin{proof}
  \uses{def:c06-item-status, def:c05-holds-value}
  A closed item's services are all delivered or removed, and neither holds
  value.
\end{proof}

\begin{theorem}[Active items are never cancelled]
  \label{thm:one-active-not-cancelled}
  \lean{Rules.OneOrder.active_not_cancelled}
  \uses{def:c06-item-status}
  An active item is not cancelled.

  Source: \href{https://standards.atlassian.net/wiki/spaces/EASD/pages/574849062/Order}{EASD guide, Order}.
\end{theorem}
\begin{proof}
  \uses{def:c06-item-status}
  The cancelled status requires that no service holds value.
\end{proof}

\begin{theorem}[Each item has exactly one status]
  \label{thm:one-item-one-status}
  \lean{Rules.OneOrder.item_has_one_status}
  \uses{def:c06-item-status, thm:one-active-not-closed, thm:one-active-not-cancelled}
  Every order item with a service has exactly one of the three statuses.

  Counterexample: an item whose one service was delivered and then
  transferred meets the conditions for both \texttt{CANCELLED} and
  \texttt{CLOSED}.  An item whose one service is \texttt{SUSPENDED} has no
  status at all under a reading where a suspended service holds no value.

  Open point: which status wins when the conditions overlap, whether a
  delivered service may be transferred at all, and what status an item has
  when none of the conditions holds.

  Source: \href{https://standards.atlassian.net/wiki/spaces/EASD/pages/574849062/Order}{EASD guide, Order}.
\end{theorem}
\begin{proof}
  \uses{def:c06-item-status}
  Not proven.  The counterexample above shows the statement is false for the
  rule as written.
\end{proof}

\begin{definition}[Order status (C-07)]
  \label{def:c07-order-status}
  \lean{Rules.OneOrder.orderStatuses}
  \uses{def:c06-item-status, def:c05-holds-value}
  An order is \texttt{OPENED} while at least one of its services holds value.
  It is \texttt{CLOSED} once all its items are closed or cancelled and all its
  services are delivered or removed.

  Source: \href{https://standards.atlassian.net/wiki/spaces/EASD/pages/574849062/Order}{EASD guide, Order}.
\end{definition}

\begin{theorem}[Opened orders are those holding value]
  \label{thm:one-order-opened}
  \lean{Rules.OneOrder.order_opened_iff}
  \uses{def:c07-order-status}
  An order is marked opened exactly when at least one of its services holds
  value.

  Source: \href{https://standards.atlassian.net/wiki/spaces/EASD/pages/574849062/Order}{EASD guide, Order}.
\end{theorem}
\begin{proof}
  \uses{def:c07-order-status}
  Unfold the status rule.
\end{proof}

\begin{theorem}[Closed orders as the guide defines them]
  \label{thm:one-order-closed}
  \lean{Rules.OneOrder.order_closed_iff}
  \uses{def:c07-order-status}
  An order is marked closed exactly when every item is closed or cancelled
  and every service is delivered or removed.

  Source: \href{https://standards.atlassian.net/wiki/spaces/EASD/pages/574849062/Order}{EASD guide, Order}.
\end{theorem}
\begin{proof}
  \uses{def:c07-order-status}
  Unfold the status rule.
\end{proof}

\begin{theorem}[New orders are opened]
  \label{thm:one-new-order-opened}
  \lean{Rules.OneOrder.accept_order_opened}
  \uses{def:one-accept, def:c07-order-status}
  Under every reading of the guide, an order created from an offer whose
  items each have a passenger and a segment is opened and not closed.

  Source: \href{https://standards.atlassian.net/wiki/spaces/EASD/pages/574849062/Order}{EASD guide, Order}.
\end{theorem}
\begin{proof}
  \uses{def:one-accept, def:c05-holds-value}
  Its services start confirmed, and a confirmed service holds value.
\end{proof}

\begin{theorem}[Orders are never both opened and closed]
  \label{thm:one-order-exclusive}
  \lean{Rules.OneOrder.opened_closed_exclusive}
  \uses{def:c07-order-status}
  Under every reading of the guide, no order is both opened and closed.

  Source: \href{https://standards.atlassian.net/wiki/spaces/EASD/pages/574849062/Order}{EASD guide, Order}.
\end{theorem}
\begin{proof}
  \uses{def:c07-order-status, def:c05-holds-value}
  An opened order has a service that holds value; a closed order's services
  are all delivered or removed, and neither holds value.
\end{proof}

\begin{theorem}[Every order is opened or closed]
  \label{thm:one-order-has-status}
  \lean{Rules.OneOrder.order_opened_or_closed}
  \uses{def:c07-order-status}
  Every order with at least one service is opened or closed.

  Counterexample: under a reading where an expired service holds no value, an
  order whose only service expired is neither opened nor closed.

  Open point: whether \texttt{EXPIRED} and \texttt{UNABLE TO DELIVER} hold
  value, and whether an order with such a service should be closed.

  Source: \href{https://standards.atlassian.net/wiki/spaces/EASD/pages/574849062/Order}{EASD guide, Order}.
\end{theorem}
\begin{proof}
  \uses{def:c07-order-status}
  Not proven.  The counterexample above shows the statement is false for the
  rule as written.
\end{proof}

% Candidate rules for offers, prices and conditions.  Written by hand; see docs/one-order-coverage.md.

\section{Offers, prices and conditions}

% Second survey round: candidates written from the survey notes in docs/one-order-survey.md.

\begin{theorem}[Offer expiry ends other limits (C-24)]
  \label{thm:offers-offer-expiry-ends-limits}
  \lean{Rules.OneOrder.Offers.guarantees_end_by_expiry}
  \uses{def:offer, def:instant, thm:accept-rejects-expired, def:offers-price-guarantee}
  Before an order exists, no price or inventory guarantee on an offer
  outlasts the offer's expiry. Once the offer's expiry instant has passed,
  both have lapsed too, unless an order was created from it. A payment time
  limit takes effect only once an order exists.

  Case to check: An offer item whose price guarantee ends an hour after the
  offer expires. A model that reads the guarantee field alone treats the
  price as guaranteed for an hour after the offer is gone.

  Open point: The Offers page also keeps the other limits alive when an
  order was changed, while the Time Limits page names only order creation.
  The business must decide whether the expiry of a reshop offer ends limits
  that would apply to the existing order it changes.

  Sources: \href{https://standards.atlassian.net/wiki/spaces/EASD/pages/600637464/Offers}{EASD guide, Offers};
  \href{https://standards.atlassian.net/wiki/spaces/EASD/pages/574423066/Time+Limits}{EASD guide, Time Limits}.
\end{theorem}
\begin{proof}
  \uses{def:offers-price-guarantee}
  A guarantee runs until the earlier of its own limit and the offer's
  expiry, so under either reading it holds at no instant after the expiry.
\end{proof}

\begin{definition}[Price guarantee at an instant (C-25)]
  \label{def:offers-price-guarantee}
  \lean{Rules.OneOrder.Offers.guaranteesAt, Rules.OneOrder.Offers.GuaranteeEnd, Rules.OneOrder.Offers.Guarantees, Rules.OneOrder.Offers.TimedOffer, Rules.OneOrder.Offers.TimedOfferItem, Rules.OneOrder.Offers.ItemLimits, Rules.OneOrder.Offers.TimedOffer.holds}
  \uses{def:offer-item, def:instant}
  An offer item's price is guaranteed at an instant only when the item
  carries a price guarantee time limit and the instant is inside both that
  limit and the offer's expiry. With no such limit the price is
  unguaranteed. The offer's expiry on its own guarantees neither price nor
  inventory. Whether a limit's own instant is inside is a reading of the
  guide, which the model takes as a parameter.

  Case to check: An item whose price guarantee ends at the same instant as
  the offer. Accepted at that instant, the offer is rejected as expired, yet
  an inclusive reading of the guarantee still calls the price guaranteed.
  The source does not fix the boundary.

  Sources: \href{https://standards.atlassian.net/wiki/spaces/EASD/pages/574423066/Time+Limits}{EASD guide, Time Limits};
  \href{https://standards.atlassian.net/wiki/spaces/EASD/pages/600637464/Offers}{EASD guide, Offers}.
\end{definition}

\begin{theorem}[No price guarantee without a limit]
  \label{thm:offers-no-limit-no-guarantee}
  \lean{Rules.OneOrder.Offers.no_limit_no_guarantee}
  \uses{def:offers-price-guarantee}
  Under every reading, an offer item with no price guarantee time limit has
  no guaranteed price at any instant.

  Source: \href{https://standards.atlassian.net/wiki/spaces/EASD/pages/574423066/Time+Limits}{EASD guide, Time Limits}.
\end{theorem}
\begin{proof}
  \uses{def:offers-price-guarantee}
  Unfold the rule: a missing limit holds nothing.
\end{proof}

\begin{theorem}[Price guaranteed before its limit]
  \label{thm:offers-price-guaranteed-before}
  \lean{Rules.OneOrder.Offers.price_guaranteed_before}
  \uses{def:offers-price-guarantee}
  Under every reading, an offer item's price is guaranteed at every instant
  before both its price guarantee time limit and the offer's expiry.

  Source: \href{https://standards.atlassian.net/wiki/spaces/EASD/pages/574423066/Time+Limits}{EASD guide, Time Limits}.
\end{theorem}
\begin{proof}
  \uses{def:offers-price-guarantee}
  An instant before both is before the earlier of the two.
\end{proof}

\begin{theorem}[Inventory guaranteed before its limit]
  \label{thm:offers-inventory-guaranteed-before}
  \lean{Rules.OneOrder.Offers.inventory_guaranteed_before}
  \uses{def:offers-price-guarantee}
  Under every reading, an offer item's inventory is guaranteed at every
  instant before both its inventory guarantee time limit and the offer's
  expiry.

  Source: \href{https://standards.atlassian.net/wiki/spaces/EASD/pages/574423066/Time+Limits}{EASD guide, Time Limits}.
\end{theorem}
\begin{proof}
  \uses{def:offers-price-guarantee}
  An instant before both is before the earlier of the two.
\end{proof}

\begin{theorem}[Guarantee ending at its instant]
  \label{thm:offers-guarantee-exclusive-acceptable}
  \lean{Rules.OneOrder.Offers.guaranteed_while_acceptable_of_exclusive}
  \uses{def:offers-price-guarantee, def:one-accept}
  Under the reading that a guarantee ends at its instant, an offer item's
  price is guaranteed only at instants at which the offer can still be
  accepted.

  Source: \href{https://standards.atlassian.net/wiki/spaces/EASD/pages/574423066/Time+Limits}{EASD guide, Time Limits}.
\end{theorem}
\begin{proof}
  \uses{def:offers-price-guarantee, def:one-accept}
  A guaranteed price means the instant is before the offer's expiry, and
  acceptance is rejected as expired only at or after it.
\end{proof}

\begin{theorem}[Guaranteed prices only on offers that can be accepted]
  \label{thm:offers-guaranteed-while-acceptable}
  \lean{Rules.OneOrder.Offers.guaranteed_while_acceptable}
  \uses{def:offers-price-guarantee, def:one-accept, thm:offers-guarantee-exclusive-acceptable}
  Under every reading, an offer item's price is guaranteed only at instants
  at which the offer can still be accepted.

  Counterexample: an offer valid until 12:00 whose flight's price guarantee
  also ends at 12:00. At 12:00 the offer is rejected as expired. Under the
  reading that a guarantee still holds at its instant, the price is
  guaranteed at that instant. The example \texttt{guaranteeAtExpiry}
  computes both.

  Open point: does a price guarantee still hold at the instant its time
  limit names?

  Source: \href{https://standards.atlassian.net/wiki/spaces/EASD/pages/574423066/Time+Limits}{EASD guide, Time Limits}.
\end{theorem}
\begin{proof}
  \uses{def:offers-price-guarantee}
  Not proven. The counterexample above shows the statement is false for the
  rule as written.
\end{proof}

\begin{definition}[Accepting after the guarantee lapses (C-26)]
  \label{def:offers-accept-after-guarantee-lapse}
  \lean{Rules.OneOrder.Offers.LapsedPrice, Rules.OneOrder.Offers.TimedOffer.guaranteeLapsed}
  \uses{def:accept, thm:accept-price-eq-sum, thm:accept-rejects-expired, def:offers-price-guarantee}
  When a customer accepts an unexpired offer after an item's price guarantee
  has lapsed, the airline decides the order price for that item. The source
  recommends repricing into a new offer before payment. The model takes the
  answer as a parameter: the acceptance is refused and sent back for a new
  offer, or the order is created at the offered price.

  Case to check: thm:accept-price-eq-sum makes the order price equal the
  offer item prices. An offer with a 30 minute price guarantee accepted two
  hours later, one minute before expiry, is sold at a price the airline no
  longer stands behind.

  Open point: The business must decide whether acceptance in this window is
  rejected, forces a reprice into a new offer, or creates the order at the
  offered price. If repricing at order creation is allowed, the order price
  can differ from the offer price.

  Sources: \href{https://standards.atlassian.net/wiki/spaces/EASD/pages/574423066/Time+Limits}{EASD guide, Time Limits};
  \href{https://standards.atlassian.net/wiki/spaces/EASD/pages/574423098/Pricing+Offers}{EASD guide, Pricing Offers}.
\end{definition}

\begin{theorem}[Order time limits only extend (C-27)]
  \label{thm:offers-order-limits-only-extend}
  \uses{def:accept, def:instant, def:c11-payment-time-limit, def:offers-price-guarantee}
  When an order is created, the time left on an offer item's price guarantee
  carries into the matching order item. After that the airline may move the
  order item's price guarantee or payment time limit later and never
  earlier.

  Case to check: An order change that recomputes the payment time limit from
  the change instant with a shorter duration than the original would move
  the limit earlier.

  Source: \href{https://standards.atlassian.net/wiki/spaces/EASD/pages/574423066/Time+Limits}{EASD guide, Time Limits}.
\end{theorem}

\begin{definition}[Payment time limit at creation (C-28)]
  \label{def:offers-payment-limit-at-creation}
  \lean{Rules.OneOrder.Offers.paymentLimitAt, Rules.OneOrder.Offers.BothForms, Rules.OneOrder.Offers.secondsAfter}
  \uses{def:accept, def:instant, def:c11-payment-time-limit, def:c12-zero-commitment}
  An offer item states its payment time limit as a duration or as a fixed
  date and time. At order creation the order item's limit becomes the
  creation instant plus the duration, or the fixed date and time. A zero
  duration means payment is due at creation. Which form applies when an
  offer item states both is a reading of the guide, which the model takes
  as a parameter.

  Case to check: A fixed offer limit that falls before the order is created
  gives an order item that is overdue the moment it exists.

  Open point: The source does not say which form wins if an offer item
  carries both, or how long an order item with no payment time limit may
  stay unpaid. The business must decide both.

  Source: \href{https://standards.atlassian.net/wiki/spaces/EASD/pages/574423066/Time+Limits}{EASD guide, Time Limits}.
\end{definition}

\begin{definition}[Time limits of a new order]
  \label{def:offers-create-order}
  \lean{Rules.OneOrder.Offers.createOrder, Rules.OneOrder.Offers.CreationReading, Rules.OneOrder.Offers.NewOrder, Rules.OneOrder.Offers.NewOrderItem, Rules.OneOrder.Offers.CreateRejection, Rules.OneOrder.Offers.newItem, Rules.OneOrder.Offers.NewOrder.toOrder, Rules.OneOrder.Offers.NewOrder.payable, Rules.OneOrder.Offers.TimedOffer.toOffer}
  \uses{def:one-accept, def:offers-price-guarantee, def:offers-accept-after-guarantee-lapse, def:offers-payment-limit-at-creation, def:c11-payment-time-limit}
  Accepting an offer whose items carry time limits runs the checks of
  ONE-001, then, under the reading that asks for a new offer, refuses the
  acceptance if an item's price guarantee has lapsed. Each order item takes
  its payment time limit and its price guarantee from its offer item.

  Source: \href{https://standards.atlassian.net/wiki/spaces/EASD/pages/574423066/Time+Limits}{EASD guide, Time Limits}.
\end{definition}

\begin{theorem}[New order matches ONE-001]
  \label{thm:offers-create-order-is-accept}
  \lean{Rules.OneOrder.Offers.createOrder_toOrder}
  \uses{def:offers-create-order, def:one-accept}
  An order created from a timed offer, without its time limits, is the order
  ONE-001 creates from the same offer, so every guarantee of ONE-001 holds
  for it.

  Source: \href{https://standards.atlassian.net/wiki/spaces/EASD/pages/574423066/Time+Limits}{EASD guide, Time Limits}.
\end{theorem}
\begin{proof}
  \uses{def:offers-create-order, def:one-accept}
  The rule creates the order only when ONE-001 does, with the same items.
\end{proof}

\begin{theorem}[Payment time limit from a duration]
  \label{thm:offers-create-duration}
  \lean{Rules.OneOrder.Offers.createOrder_duration}
  \uses{def:offers-create-order, def:offers-payment-limit-at-creation}
  An order item created from an offer item whose payment time limit is
  stated only as a duration must be paid by the creation instant plus that
  duration.

  Source: \href{https://standards.atlassian.net/wiki/spaces/EASD/pages/574423066/Time+Limits}{EASD guide, Time Limits}.
\end{theorem}
\begin{proof}
  \uses{def:offers-create-order}
  Unfold the rule for that offer item.
\end{proof}

\begin{theorem}[Payment time limit from a fixed instant]
  \label{thm:offers-create-fixed}
  \lean{Rules.OneOrder.Offers.createOrder_fixed}
  \uses{def:offers-create-order, def:offers-payment-limit-at-creation}
  An order item created from an offer item whose payment time limit is
  stated only as a fixed date and time must be paid by that instant.

  Source: \href{https://standards.atlassian.net/wiki/spaces/EASD/pages/574423066/Time+Limits}{EASD guide, Time Limits}.
\end{theorem}
\begin{proof}
  \uses{def:offers-create-order}
  Unfold the rule for that offer item.
\end{proof}

\begin{theorem}[Price guarantee carried into the order]
  \label{thm:offers-create-keeps-guarantee}
  \lean{Rules.OneOrder.Offers.createOrder_keeps_guarantee}
  \uses{def:offers-create-order}
  Each order item keeps the price guarantee time limit of its offer item.
  This is the first half of C-27; moving limits after creation is not
  modelled.

  Source: \href{https://standards.atlassian.net/wiki/spaces/EASD/pages/574423066/Time+Limits}{EASD guide, Time Limits}.
\end{theorem}
\begin{proof}
  \uses{def:offers-create-order}
  The order items are the offer items mapped one by one, each keeping its
  limit.
\end{proof}

\begin{theorem}[No order after a lapse when repricing]
  \label{thm:offers-create-reprice}
  \lean{Rules.OneOrder.Offers.createOrder_guaranteed_of_reprice}
  \uses{def:offers-create-order, def:offers-accept-after-guarantee-lapse}
  Under the reading that asks for a new offer, no order is created while
  the price guarantee of one of the offer's items has lapsed.

  Source: \href{https://standards.atlassian.net/wiki/spaces/EASD/pages/574423098/Pricing+Offers}{EASD guide, Pricing Offers}.
\end{theorem}
\begin{proof}
  \uses{def:offers-create-order}
  The rule refuses the acceptance when any item's guarantee has lapsed.
\end{proof}

\begin{theorem}[No order at a lapsed price]
  \label{thm:offers-create-price-guaranteed}
  \lean{Rules.OneOrder.Offers.createOrder_price_guaranteed}
  \uses{def:offers-create-order, thm:offers-create-reprice}
  Under every reading, no order is created at the offered price while the
  price guarantee of one of the offer's items has lapsed.

  Counterexample: an offer valid from 10:00 to 12:00 whose flight's price
  is guaranteed until 10:30, accepted at 11:59. Under the reading that
  keeps the offered price, the order is created with the flight at
  300.00 EUR. The example \texttt{lapsedGuaranteeAccepted} computes both
  readings.

  Open point: is acceptance after a lapsed price guarantee refused, sent
  back for a new offer, or honoured at the offered price?

  Sources: \href{https://standards.atlassian.net/wiki/spaces/EASD/pages/574423066/Time+Limits}{EASD guide, Time Limits};
  \href{https://standards.atlassian.net/wiki/spaces/EASD/pages/574423098/Pricing+Offers}{EASD guide, Pricing Offers}.
\end{theorem}
\begin{proof}
  \uses{def:offers-create-order}
  Not proven. The counterexample above shows the statement is false for the
  rule as written.
\end{proof}

\begin{theorem}[No order item overdue at creation]
  \label{thm:offers-create-limit-after-creation}
  \lean{Rules.OneOrder.Offers.createOrder_limit_not_before_creation}
  \uses{def:offers-create-order, def:offers-payment-limit-at-creation}
  No order item starts with a payment time limit before the order's
  creation.

  Counterexample: an offer valid until 12:00 whose bag must be paid by
  10:30, accepted at 11:00. The bag's order item must be paid half an hour
  before the order existed. The example \texttt{fixedLimitBeforeCreation}
  computes it.

  Open point: is an acceptance after a fixed payment time limit refused, or
  does the order item get a new payment time limit?

  Source: \href{https://standards.atlassian.net/wiki/spaces/EASD/pages/574423066/Time+Limits}{EASD guide, Time Limits}.
\end{theorem}
\begin{proof}
  \uses{def:offers-create-order}
  Not proven. The counterexample above shows the statement is false for the
  rule as written.
\end{proof}

\begin{theorem}[Payment due at creation can be made]
  \label{thm:offers-zero-duration-payable}
  \lean{Rules.OneOrder.Offers.zero_duration_payable}
  \uses{def:offers-create-order, def:offers-payment-limit-at-creation, def:payment-accept}
  An order item whose offer item's payment time limit is a duration of zero
  can be paid at the instant the order is created.

  Counterexample: a flight at 300.00 EUR with a payment time limit of zero,
  accepted at 11:00 and paid in full by card at 11:00. The order item's
  limit is 11:00, and the payment rule of ONE-304 rejects a payment made at
  or after the limit. The example \texttt{zeroDurationPayment} computes it.

  Open point: is a payment made at the creation instant on time, or must it
  come with the request that creates the order?

  Source: \href{https://standards.atlassian.net/wiki/spaces/EASD/pages/574423066/Time+Limits}{EASD guide, Time Limits}.
\end{theorem}
\begin{proof}
  \uses{def:offers-create-order, def:payment-accept}
  Not proven. The counterexample above shows the statement is false for the
  rule as written.
\end{proof}

\begin{theorem}[Nesting of guarantee time limits (C-29)]
  \label{thm:offers-guarantee-limits-nest}
  \uses{def:offers-price-guarantee, def:c11-payment-time-limit}
  \notready
  When an item has a price guarantee time limit, that limit is at or after
  every separate guarantee limit on its base amount, each tax and each fee.
  Those part limits count only when the overall limit is present. The price
  guarantee should also end no later than the payment time limit.

  Case to check: The source's table guarantees the base amount and one tax
  for 5 days in the order. If the order item's overall guarantee stays at
  the 30 minutes carried from the offer, the parts outlive the whole.

  Open point: The source calls the overall limit the maximum of the parts,
  which reads either as equal to the latest part or as an upper bound, and
  the payment limit bound is only a recommendation. The business must pick
  one reading and decide whether a price guarantee past the payment limit is
  cut back or refused.

  Sources: \href{https://standards.atlassian.net/wiki/spaces/EASD/pages/574423066/Time+Limits}{EASD guide, Time Limits};
  \href{https://standards.atlassian.net/wiki/spaces/EASD/pages/574914636/Fees}{EASD guide, Fees}.
\end{theorem}

\begin{definition}[Restriction rule in force (C-30)]
  \label{def:offers-restriction-in-force}
  \lean{Rules.OneOrder.Offers.ruleInForce, Rules.OneOrder.Offers.Restriction, Rules.OneOrder.Offers.Stage, Rules.OneOrder.Offers.Journey, Rules.OneOrder.Offers.Journey.stageAt, Rules.OneOrder.Offers.Action, Rules.OneOrder.Offers.ChangeKind, Rules.OneOrder.Offers.ExpiryReading, Rules.OneOrder.Offers.Restriction.live, Rules.OneOrder.Offers.latestOf, Rules.OneOrder.Offers.candidates}
  \uses{def:offer-item, def:instant}
  Each change or cancel rule on an offer item belongs to one journey stage
  (before payment, before departure, after departure, no show) and starts at
  its effective date. At a given instant the rule in force is the one in the
  current stage with the latest effective date not after that instant. A
  rule stops when its stage ends or a later rule takes over. Whether a
  rule's expiration date also ends it is a reading of the guide, which the
  model takes as a parameter.

  Case to check: In the source example (payment 1 Jan 2019, departure 1 Feb
  2019) the no-change rule described as 7 days before departure starts on 24
  Jan, which is 8 days before. The text gives the middle fee as 1 CHF where
  the table and diagram give 60 CHF. Another example sets equal effective
  and expiration dates, so that rule is never in force.

  Open point: The message examples also give rules an expiration date, which
  the stage and override reading never uses. The business must decide which
  governs, and what applies before the first effective date or in a gap
  between one rule's expiration and the next rule. The source also does not
  say which of two rules of a stage applies when they start at the same
  instant.

  Sources: \href{https://standards.atlassian.net/wiki/spaces/EASD/pages/574259692/Offer+Restrictions}{EASD guide, Offer Restrictions};
  \href{https://standards.atlassian.net/wiki/spaces/EASD/pages/574816307/Offer+Conditions+and+Restrictions}{EASD guide, Offer Conditions and Restrictions}.
\end{definition}

\begin{theorem}[Rule in force belongs to the current stage]
  \label{thm:offers-rule-stage}
  \lean{Rules.OneOrder.Offers.ruleInForce_stage}
  \uses{def:offers-restriction-in-force}
  The rule in force for an action at an instant is one of the rules for
  that action and for the stage the journey is in at that instant.

  Source: \href{https://standards.atlassian.net/wiki/spaces/EASD/pages/574259692/Offer+Restrictions}{EASD guide, Offer Restrictions}.
\end{theorem}
\begin{proof}
  \uses{def:offers-restriction-in-force}
  The rule picks from the rules filtered by action and stage.
\end{proof}

\begin{theorem}[Rule in force has started]
  \label{thm:offers-rule-started}
  \lean{Rules.OneOrder.Offers.ruleInForce_started}
  \uses{def:offers-restriction-in-force}
  The rule in force at an instant has its effective date at or before that
  instant.

  Source: \href{https://standards.atlassian.net/wiki/spaces/EASD/pages/574259692/Offer+Restrictions}{EASD guide, Offer Restrictions}.
\end{theorem}
\begin{proof}
  \uses{def:offers-restriction-in-force}
  The rule picks only from rules that have started.
\end{proof}

\begin{theorem}[Latest started rule is in force]
  \label{thm:offers-rule-latest}
  \lean{Rules.OneOrder.Offers.ruleInForce_latest}
  \uses{def:offers-restriction-in-force}
  When rules for an action and the current stage have started, and have not
  expired where expiration dates end rules, one of them is in force, and
  none of them starts later than it.

  Source: \href{https://standards.atlassian.net/wiki/spaces/EASD/pages/574259692/Offer+Restrictions}{EASD guide, Offer Restrictions}.
\end{theorem}
\begin{proof}
  \uses{def:offers-restriction-in-force}
  By induction on the list of rules: the latest of a list is one of them and
  starts no earlier than any of them.
\end{proof}

\begin{theorem}[Expired rules are not in force]
  \label{thm:offers-rule-not-expired}
  \lean{Rules.OneOrder.Offers.ruleInForce_not_expired}
  \uses{def:offers-restriction-in-force}
  Under the reading that expiration dates end rules, the rule in force at an
  instant has not reached its expiration date.

  Source: \href{https://standards.atlassian.net/wiki/spaces/EASD/pages/574816307/Offer+Conditions+and+Restrictions}{EASD guide, Offer Conditions and Restrictions}.
\end{theorem}
\begin{proof}
  \uses{def:offers-restriction-in-force}
  Under that reading the rule picks only from rules before their expiration
  date.
\end{proof}

\begin{theorem}[Every stage with rules has a rule in force]
  \label{thm:offers-rule-covers-stage}
  \lean{Rules.OneOrder.Offers.ruleInForce_covers_stage}
  \uses{def:offers-restriction-in-force, thm:offers-rule-latest}
  While a journey is in a stage for which an action has rules, a rule is in
  force for that action.

  Counterexample: a journey paid at 10:00 on 1 January whose only flight
  change rule before departure starts on 2 January has no rule in force at
  18:00 on 1 January, under either reading. Under the reading that
  expiration dates end rules, a rule whose effective and expiration dates
  are equal is never in force. The example \texttt{noRuleInForce} computes
  both.

  Open point: what applies to a change or cancellation when no rule of the
  current stage is in force?

  Sources: \href{https://standards.atlassian.net/wiki/spaces/EASD/pages/574259692/Offer+Restrictions}{EASD guide, Offer Restrictions};
  \href{https://standards.atlassian.net/wiki/spaces/EASD/pages/574816307/Offer+Conditions+and+Restrictions}{EASD guide, Offer Conditions and Restrictions}.
\end{theorem}
\begin{proof}
  \uses{def:offers-restriction-in-force}
  Not proven. The counterexample above shows the statement is false for the
  rule as written.
\end{proof}

\begin{definition}[Voluntary change under the rule in force]
  \label{def:offers-voluntary-change}
  \lean{Rules.OneOrder.Offers.voluntaryChange, Rules.OneOrder.Offers.ChangeRequest, Rules.OneOrder.Offers.ChangeDecision, Rules.OneOrder.Offers.Restriction.feeOr}
  \uses{def:offers-restriction-in-force, def:change-cancellation-price-differential}
  A voluntary change is refused when the rule in force forbids it. When the
  rule allows it, the change goes to the airline's pricing of ONE-201 with
  the old and new prices, the rule's fee as the penalty and nothing
  forfeited. When no rule is in force, the decision says so.

  Source: \href{https://standards.atlassian.net/wiki/spaces/EASD/pages/574259692/Offer+Restrictions}{EASD guide, Offer Restrictions}.
\end{definition}

\begin{theorem}[Forbidden changes are refused]
  \label{thm:offers-change-refused}
  \lean{Rules.OneOrder.Offers.change_refused}
  \uses{def:offers-voluntary-change}
  A change the rule in force forbids is refused.

  Source: \href{https://standards.atlassian.net/wiki/spaces/EASD/pages/574259692/Offer+Restrictions}{EASD guide, Offer Restrictions}.
\end{theorem}
\begin{proof}
  \uses{def:offers-voluntary-change}
  Unfold the rule.
\end{proof}

\begin{theorem}[Rule's fee is the penalty]
  \label{thm:offers-change-fee}
  \lean{Rules.OneOrder.Offers.change_fee_is_penalty}
  \uses{def:offers-voluntary-change}
  A change the rule in force allows with a fee is quoted with that fee as
  its penalty.

  Source: \href{https://standards.atlassian.net/wiki/spaces/EASD/pages/574259692/Offer+Restrictions}{EASD guide, Offer Restrictions}.
\end{theorem}
\begin{proof}
  \uses{def:offers-voluntary-change}
  Unfold the rule.
\end{proof}

\begin{theorem}[No fee, no penalty]
  \label{thm:offers-change-no-fee}
  \lean{Rules.OneOrder.Offers.change_no_fee}
  \uses{def:offers-voluntary-change}
  A change the rule in force allows without a fee is quoted with no
  penalty.

  Source: \href{https://standards.atlassian.net/wiki/spaces/EASD/pages/574259692/Offer+Restrictions}{EASD guide, Offer Restrictions}.
\end{theorem}
\begin{proof}
  \uses{def:offers-voluntary-change}
  Unfold the rule.
\end{proof}

\begin{theorem}[Voluntary change under a restriction (C-31)]
  \label{thm:offers-change-under-restriction}
  \lean{Rules.OneOrder.Offers.change_net_is_fee_plus_difference}
  \uses{def:offers-restriction-in-force, def:offers-voluntary-change, def:c18-change, def:change-cancellation-price-differential, def:change-cancellation-due-amounts}
  A voluntary change of a given type, such as flight or name, is rejected
  when the rule in force forbids it. When the rule allows it, the customer
  pays the rule's fee, or nothing if the rule states no fee, plus any fare
  difference. The model proves the payment for a fee no larger than the
  price of the items changed, when the seller does not ask to keep the value
  for later use.

  Case to check: The source example described as a no-show change with total
  loss of fare carries a name change rule with a 1 CHF fee and no stage, so
  its description and its data disagree on what the customer owes.

  Open point: The source does not say what applies when no rule covers the
  current stage, or whether a fee marked approximate binds when the change
  is made. The business must choose a default for an uncovered stage and
  decide how an approximate fee is settled.

  Sources: \href{https://standards.atlassian.net/wiki/spaces/EASD/pages/574259692/Offer+Restrictions}{EASD guide, Offer Restrictions};
  \href{https://standards.atlassian.net/wiki/spaces/EASD/pages/574816307/Offer+Conditions+and+Restrictions}{EASD guide, Offer Conditions and Restrictions}.
\end{theorem}
\begin{proof}
  \uses{def:offers-voluntary-change, def:change-cancellation-due-amounts}
  The quote's price difference is the new price less the old, and a fee no
  larger than the old price is charged in full under both readings of
  ONE-202, netted or separate.
\end{proof}

\begin{definition}[Refund owed on voluntary cancellation (C-32)]
  \label{def:offers-refund-on-cancellation}
  \lean{Rules.OneOrder.Offers.cancellationRefund, Rules.OneOrder.Offers.RefundTerms, Rules.OneOrder.Offers.RefundTerms.amount, Rules.OneOrder.Offers.FullOrFixed, Rules.OneOrder.Offers.RefundDecision, Rules.OneOrder.Offers.CancelledItem, Rules.OneOrder.Offers.CancelledItem.refundableCharges}
  \uses{def:offers-restriction-in-force, def:c15-cancel-order, thm:c16-penalty-net, def:refund}
  The cancel rule in force decides what a voluntary cancellation returns. A
  rule that forbids cancellation returns nothing. A rule that allows it
  returns a full or a fixed refund, paid to the original form of payment
  when the rule says so. Which amount a refund marked full and also given a
  fixed amount returns is a reading of the guide, which the model takes as
  a parameter.

  Case to check: Reading that example as the full item price and reading it
  as 1.00 CHF give two different refunds for the same rule.

  Open point: One source example marks a refund as full and also gives a
  fixed amount of 1.00 CHF. The business must decide which governs when both
  are present, and how the rule's refund relates to the penalty netting in
  thm:c16-penalty-net.

  Source: \href{https://standards.atlassian.net/wiki/spaces/EASD/pages/574259692/Offer+Restrictions}{EASD guide, Offer Restrictions}.
\end{definition}

\begin{theorem}[Forbidden cancellation returns nothing]
  \label{thm:offers-cancel-forbidden}
  \lean{Rules.OneOrder.Offers.cancel_forbidden_nothing}
  \uses{def:offers-refund-on-cancellation}
  A cancellation the rule in force forbids returns nothing.

  Source: \href{https://standards.atlassian.net/wiki/spaces/EASD/pages/574259692/Offer+Restrictions}{EASD guide, Offer Restrictions}.
\end{theorem}
\begin{proof}
  \uses{def:offers-refund-on-cancellation}
  Unfold the rule.
\end{proof}

\begin{theorem}[Full refund returns the whole price]
  \label{thm:offers-cancel-full}
  \lean{Rules.OneOrder.Offers.cancel_full_refund}
  \uses{def:offers-refund-on-cancellation}
  A cancellation the rule in force allows with a full refund and no fixed
  amount returns the item's whole price.

  Source: \href{https://standards.atlassian.net/wiki/spaces/EASD/pages/574259692/Offer+Restrictions}{EASD guide, Offer Restrictions}.
\end{theorem}
\begin{proof}
  \uses{def:offers-refund-on-cancellation}
  Unfold the rule.
\end{proof}

\begin{theorem}[Fixed refund returns its amount]
  \label{thm:offers-cancel-fixed}
  \lean{Rules.OneOrder.Offers.cancel_fixed_refund}
  \uses{def:offers-refund-on-cancellation}
  A cancellation the rule in force allows with a fixed refund not marked
  full returns that amount.

  Source: \href{https://standards.atlassian.net/wiki/spaces/EASD/pages/574259692/Offer+Restrictions}{EASD guide, Offer Restrictions}.
\end{theorem}
\begin{proof}
  \uses{def:offers-refund-on-cancellation}
  Unfold the rule.
\end{proof}

\begin{theorem}[Refund to the original form of payment]
  \label{thm:offers-cancel-original-form}
  \lean{Rules.OneOrder.Offers.cancel_to_original_form}
  \uses{def:offers-refund-on-cancellation}
  A cancellation the rule in force allows, with the refund sent to the
  original form of payment, is refunded to the original form of payment.

  Source: \href{https://standards.atlassian.net/wiki/spaces/EASD/pages/574259692/Offer+Restrictions}{EASD guide, Offer Restrictions}.
\end{theorem}
\begin{proof}
  \uses{def:offers-refund-on-cancellation}
  Unfold the rule.
\end{proof}

\begin{theorem}[Taxes on a non-refundable fare (C-33)]
  \label{thm:offers-refundable-taxes}
  \lean{Rules.OneOrder.Offers.cancel_returns_refundable_charges}
  \uses{def:offers-refund-on-cancellation, thm:offers-tax-summary-total, def:refund}
  Each tax, fee and charge on a price carries its own refund flag. When an
  unused item is cancelled under a rule that forfeits the fare, the customer
  is still owed the taxes, fees and charges flagged refundable.

  Counterexample: a trip at 500.00 EUR with a refundable airport tax of
  25.00 EUR, cancelled on 25 January under a rule that forbids
  cancellation, returns nothing; so does a cancellation after a no show.
  Under the reading that a refund marked full with a fixed amount of
  1.00 EUR returns the fixed amount, the cancellation returns 1.00 EUR. The
  examples \texttt{nonRefundableFareTaxes} and \texttt{fullMarkedWithAmount}
  compute them.

  Open point: The source gives refund flags per tax and per fee and refund
  rules per offer item, but never says how they combine. The business must
  decide whether a no-refund cancel rule also keeps refundable taxes, and
  whether a no show changes the answer.

  Sources: \href{https://standards.atlassian.net/wiki/spaces/EASD/pages/574455915/Taxes+Fees+and+Charges}{EASD guide, Taxes, Fees, and Charges};
  \href{https://standards.atlassian.net/wiki/spaces/EASD/pages/574914636/Fees}{EASD guide, Fees};
  \href{https://standards.atlassian.net/wiki/spaces/EASD/pages/574259692/Offer+Restrictions}{EASD guide, Offer Restrictions}.
\end{theorem}
\begin{proof}
  \uses{def:offers-refund-on-cancellation}
  Not proven. The counterexample above shows the statement is false for the
  rule as written.
\end{proof}

\begin{definition}[Totalling a price from its parts]
  \label{def:offers-price-parts}
  \lean{Rules.OneOrder.Offers.priceTotal, Rules.OneOrder.Offers.PriceParts, Rules.OneOrder.Offers.PriceTotal, Rules.OneOrder.Offers.PriceTotal.of, Rules.OneOrder.Offers.PriceRejection, Rules.OneOrder.Offers.Charge, Rules.OneOrder.Offers.ConsumptionTax, Rules.OneOrder.Offers.ExcludedTax, Rules.OneOrder.Offers.ExcludedTax.amount, Rules.OneOrder.Offers.PriceParts.lineExcluded, Rules.OneOrder.Offers.PriceParts.amounts, Rules.OneOrder.Offers.sumOf}
  \uses{def:money, def:money-add, def:offer-item}
  A price stated with its parts is rejected when its amounts are not all in
  one currency or its discount is larger than the amount before discounts.
  Otherwise its base amount is the amount before discounts less the
  discount, its tax summary total is the sum of its entries, and its total
  adds the base amount, the surcharges, the fees, the tax summary total and
  the consumption tax not included in it. Whether that tax comes from the
  lines marked as not included or from the stated total of excluded tax is
  a reading of the guide, which the model takes as a parameter. Discounts
  are stated as amounts.

  Sources: \href{https://standards.atlassian.net/wiki/spaces/EASD/pages/574455915/Taxes+Fees+and+Charges}{EASD guide, Taxes, Fees, and Charges};
  \href{https://standards.atlassian.net/wiki/spaces/EASD/pages/2588377182/Consumption+Taxes}{EASD guide, Consumption Taxes}.
\end{definition}

\begin{theorem}[Mixed-currency prices are rejected]
  \label{thm:offers-price-mixed-currency}
  \lean{Rules.OneOrder.Offers.price_rejects_mixed_currency}
  \uses{def:offers-price-parts}
  A price with an amount in a currency other than that of its amount before
  discounts is rejected.

  Source: \href{https://standards.atlassian.net/wiki/spaces/EASD/pages/749305864/Currency+in+Monetary+Amounts}{EASD guide, Currency in Monetary Amounts}.
\end{theorem}
\begin{proof}
  \uses{def:offers-price-parts}
  The currency check runs first.
\end{proof}

\begin{theorem}[Discount above the base is rejected]
  \label{thm:offers-price-large-discount}
  \lean{Rules.OneOrder.Offers.price_rejects_large_discount}
  \uses{def:offers-price-parts}
  A price whose discount is larger than its amount before discounts is never
  totalled.

  Source: \href{https://standards.atlassian.net/wiki/spaces/EASD/pages/574849038/Discounts}{EASD guide, Discounts}.
\end{theorem}
\begin{proof}
  \uses{def:offers-price-parts}
  Every branch the rule can reach for such a price is a rejection.
\end{proof}

\begin{theorem}[Base amount after discounts]
  \label{thm:offers-price-base}
  \lean{Rules.OneOrder.Offers.price_base}
  \uses{def:offers-price-parts}
  A totalled price's base amount is its amount before discounts less its
  discount.

  Source: \href{https://standards.atlassian.net/wiki/spaces/EASD/pages/574849038/Discounts}{EASD guide, Discounts}.
\end{theorem}
\begin{proof}
  \uses{def:offers-price-parts}
  Unfold the rule.
\end{proof}

\begin{theorem}[Tax summary total from its entries]
  \label{thm:offers-price-tax-total}
  \lean{Rules.OneOrder.Offers.price_tax_total}
  \uses{def:offers-price-parts}
  A totalled price's tax summary total is the sum of its entries.

  Source: \href{https://standards.atlassian.net/wiki/spaces/EASD/pages/574455915/Taxes+Fees+and+Charges}{EASD guide, Taxes, Fees, and Charges}.
\end{theorem}
\begin{proof}
  \uses{def:offers-price-parts}
  Unfold the rule.
\end{proof}

\begin{theorem}[Total from base, surcharges, fees and taxes]
  \label{thm:offers-price-total-adds}
  \lean{Rules.OneOrder.Offers.price_total_from_parts}
  \uses{def:offers-price-parts}
  A totalled price's total is its base amount plus its surcharges, its
  fees, its tax summary total and the consumption tax not included in it,
  as the reading takes that tax, all in the price's currency.

  Source: \href{https://standards.atlassian.net/wiki/spaces/EASD/pages/574455915/Taxes+Fees+and+Charges}{EASD guide, Taxes, Fees, and Charges}.
\end{theorem}
\begin{proof}
  \uses{def:offers-price-parts}
  Unfold the rule.
\end{proof}

\begin{theorem}[Excluded tax from the lines]
  \label{thm:offers-price-lines-reading}
  \lean{Rules.OneOrder.Offers.price_adds_excluded_lines_of_lines}
  \uses{def:offers-price-parts}
  Under the reading that takes the excluded consumption tax from the lines,
  a totalled price adds the tax of every line marked as not included.

  Source: \href{https://standards.atlassian.net/wiki/spaces/EASD/pages/2588377182/Consumption+Taxes}{EASD guide, Consumption Taxes}.
\end{theorem}
\begin{proof}
  \uses{def:offers-price-parts}
  Unfold the rule under that reading.
\end{proof}

\begin{theorem}[Price total from its parts (C-34)]
  \label{thm:offers-price-total-from-parts}
  \lean{Rules.OneOrder.Offers.price_adds_excluded_lines}
  \uses{def:money-add, def:offer-item, def:offers-price-parts, thm:offers-price-total-adds, thm:offers-price-lines-reading, thm:offers-tax-summary-total, thm:offers-consumption-tax-lines}
  A price's total equals its base amount plus surcharges, fees, the tax
  summary total and any consumption tax not already included. The base
  amount is stated after discounts: the pre-discount amount minus the
  discount, where a percentage discount is that share of the pre-discount
  amount.

  Case to check: In the Canadian example every HST line is marked as not
  included, yet each reports zero excluded tax. Adding the excluded tax
  raises the total by 58.78 CAD; trusting the zero totals leaves it
  unchanged.

  Counterexample: a base of 452.15 with an HST line of 58.78 marked as not
  included and a stated total of excluded tax of zero, in US dollars since
  the model has no Canadian dollar. Taking the tax from the lines gives a
  total of 510.93; trusting the stated total gives 452.15. The example
  \texttt{excludedTaxReadings} computes both.

  Open point: The source names the parts and gives one example that adds up
  (95.00 base and 20.00 tax give 115.00), but states no formula. The
  business must decide whether fees and surcharges sit inside the tax
  summary or add separately, whether a discount applies before taxes are
  computed, and whether excluded consumption tax adds to the total. The
  model states discounts as amounts, so a percentage discount and its
  rounding are not modelled.

  Sources: \href{https://standards.atlassian.net/wiki/spaces/EASD/pages/574455915/Taxes+Fees+and+Charges}{EASD guide, Taxes, Fees, and Charges};
  \href{https://standards.atlassian.net/wiki/spaces/EASD/pages/574914636/Fees}{EASD guide, Fees};
  \href{https://standards.atlassian.net/wiki/spaces/EASD/pages/574849038/Discounts}{EASD guide, Discounts};
  \href{https://standards.atlassian.net/wiki/spaces/EASD/pages/2588377182/Consumption+Taxes}{EASD guide, Consumption Taxes}.
\end{theorem}
\begin{proof}
  \uses{def:offers-price-parts}
  Not proven. The counterexample above shows the statement is false for the
  rule as written.
\end{proof}

\begin{theorem}[Tax summary total and detail (C-35)]
  \label{thm:offers-tax-summary-total}
  \uses{def:money-add, def:accept, def:order-item}
  A tax summary's total equals the sum of its individual tax, fee and charge
  amounts, all in one currency. A shopping response may give only the total,
  but the order created from an accepted offer must hold the individual
  entries, including consumption tax detail.

  Case to check: An offer shopped with only a tax total and accepted without
  a pricing step has no entries to copy into the order. The source does not
  say where they come from or require them to add up to the shopped total.

  Sources: \href{https://standards.atlassian.net/wiki/spaces/EASD/pages/574455915/Taxes+Fees+and+Charges}{EASD guide, Taxes, Fees, and Charges};
  \href{https://standards.atlassian.net/wiki/spaces/EASD/pages/2588377182/Consumption+Taxes}{EASD guide, Consumption Taxes}.
\end{theorem}

\begin{theorem}[Consumption tax lines add up (C-36)]
  \label{thm:offers-consumption-tax-lines}
  \uses{def:money, def:money-add, thm:offers-tax-summary-total}
  \notready
  In each consumption tax entry, the tax is the rate applied to the taxable
  base, and the entry's total tax is the sum of its line taxes. A gross
  amount equals its net amount plus every consumption tax charged on it. All
  amounts in an entry are in the price's currency.

  Case to check: Italian example: line tax 1.82 against entry total 1.80.
  Canadian example: line tax 0.93 against entry total 4.55. Australian
  example: gross in AUD with tax in EUR. Tahiti example: two taxes on one
  base, so a per-line check that gross equals net plus that line's tax
  fails.

  Open point: The source shows these relations only through examples,
  several of which break them, and gives no rounding rule. The business must
  decide which figure is authoritative when a line and its total disagree,
  and how to round.

  Source: \href{https://standards.atlassian.net/wiki/spaces/EASD/pages/2588377182/Consumption+Taxes}{EASD guide, Consumption Taxes}.
\end{theorem}

\begin{theorem}[Seller margin stays outside order (C-37)]
  \label{thm:offers-margin-outside-order}
  \uses{def:accept, thm:accept-price-eq-sum, def:payment}
  When the seller is merchant of record and adds a margin, the airline's
  order total equals the airline offer price including taxes, fees and
  charges. The order holds no record of the margin, and the seller pays the
  airline the offer price.

  Case to check: The source places the airline as merchant of record out of
  scope. In that case the payment recorded on the order would include the
  margin and exceed the order total, which a rule tying payments to the
  order total would reject.

  Source: \href{https://standards.atlassian.net/wiki/spaces/EASD/pages/574849050/Commission+and+Margin}{EASD guide, Commission and Margin}.
\end{theorem}

\begin{definition}[Mandatory and optional offer items (C-38)]
  \label{def:offers-mandatory-optional-items}
  \uses{def:offer, def:offer-item, def:accept, thm:c02-item-per-offer-item, thm:accept-price-eq-sum, thm:c04-item-unsplit}
  \notready
  Each offer item is mandatory or optional, and its services and price are
  fixed once offered. Accepting an offer takes every mandatory item whole
  plus the optional items chosen; a-la-carte items are all optional, priced
  per unit, and bought in a chosen quantity. Every offer other than an
  a-la-carte offer should hold at least one mandatory item.

  Case to check: thm:c02-item-per-offer-item creates one order item per
  offer item, which fails when the customer skips an optional item. The
  source's own figure shows an offer total of 80 USD equal to one item while
  a second item is listed.

  Open point: The source does not say whether an offer's total price counts
  its optional items. The business must decide that, and how acceptance
  records the chosen items and quantities.

  Sources: \href{https://standards.atlassian.net/wiki/spaces/EASD/pages/600637464/Offers}{EASD guide, Offers};
  \href{https://standards.atlassian.net/wiki/spaces/SRSIA/pages/574619905/General+Principles+and+Assumptions}{Interline guide, General Principles and Assumptions}.
\end{definition}

\begin{theorem}[Offer responsible airline owns order (C-39)]
  \label{thm:offers-ora-owns-order}
  \uses{def:offer, def:accept, def:order}
  Each offer has one Offer Responsible Airline, the airline that returns the
  combined offer; every other airline in it is a Participating Offer
  Airline. The order created from the offer is owned by the Offer
  Responsible Airline, which receives the seller's payment commitment and
  holds the single source of the order's data.

  Case to check: In the interline ordering flow each participating airline
  also keeps its own order for its services. Status read from that order can
  disagree with the owner's order unless every status change reaches the
  owner.

  Sources: \href{https://www.iata.org/contentassets/72cbd60393ff42b5975d90ce9e049a7d/oneorder-resolution-797.pdf}{IATA Resolution 797};
  \href{https://standards.atlassian.net/wiki/spaces/SWO/pages/574816490/Roles}{Settlement with Orders, Roles};
  \href{https://standards.atlassian.net/wiki/spaces/SRSIA/pages/574685353/Ordering}{Interline guide, Ordering}.
\end{theorem}

\begin{definition}[Terms of a combined interline offer (C-40)]
  \label{def:offers-interline-offer-terms}
  \uses{thm:offers-ora-owns-order, def:offers-restriction-in-force, def:offers-mandatory-optional-items}
  \notready
  When the Offer Responsible Airline combines a participating airline's
  offer item into its own offer, it keeps that item whole and accepts the
  item's settlement value, settlement terms and sales conditions. It sets
  the customer price itself, and the settlement value, held per service,
  stays between the two airlines.

  Case to check: A participating airline's segment is non-refundable while
  the combined offer item is shown as fully refundable. The customer gets a
  full refund and the Offer Responsible Airline still owes the settlement
  value.

  Open point: The source does not say how the participating airline's change
  and cancel rules become the restrictions shown to the customer, or whether
  the customer price must cover the settlement values. The business must
  decide both.

  Sources: \href{https://standards.atlassian.net/wiki/spaces/SRSIA/pages/574816454/Shopping}{Interline guide, Shopping};
  \href{https://standards.atlassian.net/wiki/spaces/SRSIA/pages/574619905/General+Principles+and+Assumptions}{Interline guide, General Principles and Assumptions};
  \href{https://standards.atlassian.net/wiki/spaces/SRSIA/pages/574914728/Glossary+of+Terms}{Interline guide, Glossary of Terms};
  \href{https://standards.atlassian.net/wiki/spaces/SRSIA/overview}{Interline guide}.
\end{definition}

\begin{definition}[Currency of an offer's amounts (C-41)]
  \label{def:offers-offer-currency}
  \uses{def:money, def:currency, def:offer, def:money-add}
  \notready
  Every amount carries a currency, zero included. Amounts that add into an
  offer's totals are in the offer's currency and use that currency's number
  of decimals (two for GBP, three for TND, none for JPY). Local or
  equivalent amounts in another currency stay out of totals.

  Case to check: The source's Tahiti example states XPF amounts to two
  decimals, although ISO 4217 gives XPF none. A minor-unit model reads
  100000.00 XPF as either 100000 or 10000000 units depending on which rule
  it trusts.

  Open point: The source lets a seller request a currency and lets a fee
  carry a local amount, but never says which amount binds or whether a
  conversion rate is guaranteed. The business must decide what the customer
  pays when the requested currency differs from the airline's pricing
  currency.

  Sources: \href{https://standards.atlassian.net/wiki/spaces/EASD/pages/749305864/Currency+in+Monetary+Amounts}{EASD guide, Currency in Monetary Amounts};
  \href{https://standards.atlassian.net/wiki/spaces/EASD/pages/574685264/Price+Decimals}{EASD guide, Price Decimals};
  \href{https://standards.atlassian.net/wiki/spaces/EASD/pages/574914636/Fees}{EASD guide, Fees};
  \href{https://retailing.iata.org/armi/docs/SHPLOC/}{IATA capability SHPLOC}.
\end{definition}

\begin{definition}[Time limits as UTC instants (C-42)]
  \label{def:offers-time-limits-as-instants}
  \uses{def:instant, thm:offers-offer-expiry-ends-limits, def:offers-payment-limit-at-creation}
  \notready
  Every time limit on an offer or order resolves to one UTC instant before
  it is compared with the current time. Flight times are local to their
  airport and need its offset to resolve.

  Case to check: A payment limit of 2018-01-01 read as midnight UTC and read
  as the end of that day in a seller zone of UTC-12 differ by 36 hours, so
  one payment is on time under one reading and late under the other.

  Open point: The source prefers UTC, but its own time limit examples
  include a bare date and a date and time with no offset. The business must
  decide which instant such a limit denotes, for example the end of that day
  at the point of sale.

  Sources: \href{https://standards.atlassian.net/wiki/spaces/EASD/pages/574488607/Reference+to+Times}{EASD guide, Reference to Times};
  \href{https://standards.atlassian.net/wiki/spaces/EASD/pages/574423066/Time+Limits}{EASD guide, Time Limits}.
\end{definition}

% Candidate rules for ancillaries, seats and bundles.  Written by hand; see docs/one-order-coverage.md.

\section{Ancillaries, seats and bundles}

% No candidates yet.

% Candidate rules for order creation and information.  Written by hand; see docs/one-order-coverage.md.

\section{Order creation and information}

\begin{theorem}[Order version (C-21)]
  \label{thm:c21-version}
  \uses{def:c18-change, def:order}
  An order's version starts at 1 and goes up by one for each group of changes,
  so the version is always one more than the number of change groups applied.

  Source: \href{https://standards.atlassian.net/wiki/spaces/EASD/pages/574259659/Order+Versioning}{EASD guide, Order Versioning}.
\end{theorem}

\begin{definition}[Order retention (C-23)]
  \label{def:c23-retention}
  \uses{def:c07-order-status, def:instant}
  An order with an open item stays accessible for at least 13 months from
  creation.  Once all its items reach a final status, it stays accessible for
  at least 7 more days.

  Source: \href{https://www.iata.org/contentassets/72cbd60393ff42b5975d90ce9e049a7d/oneorder-resolution-797.pdf}{IATA Resolution 797}.
\end{definition}

% Candidate rules for change and cancellation.  Written by hand; see docs/one-order-coverage.md.

\section{Change and cancellation}

\begin{definition}[Cancelling an order (C-15)]
  \label{def:c15-cancel-order}
  \lean{Rules.OneOrder.Change.requestCancellation, Rules.OneOrder.Change.cancelDirect, Rules.OneOrder.Change.CancelTerms, Rules.OneOrder.Change.cancellationQuote}
  \uses{def:c06-item-status, def:payment, def:totals}
  An unpaid order can be cancelled directly.  A paid order is cancelled only by
  accepting a cancellation offer from the airline.  The model cancels an order
  at once when no payment has been received for it, cancelling and removing
  every service.  Otherwise it prices a cancellation offer from the order's
  price, as a change whose new price is zero.

  Source: \href{https://standards.atlassian.net/wiki/spaces/EASD/pages/574783569/Cancellation+of+Order+Items}{EASD guide, Cancellation of Order Items}.

  Capabilities: ORDCAN, ORDCA2.
\end{definition}

\begin{theorem}[Unpaid orders are cancelled at once]
  \label{thm:change-cancellation-unpaid-at-once}
  \lean{Rules.OneOrder.Change.unpaid_cancelled_at_once}
  \uses{def:c15-cancel-order}
  When no payment has been received for an order, a request to cancel it
  cancels and removes every service at once, and the order keeps its items.

  Source: \href{https://standards.atlassian.net/wiki/spaces/EASD/pages/574783569/Cancellation+of+Order+Items}{EASD guide, Cancellation of Order Items}.
\end{theorem}
\begin{proof}
  \uses{def:c15-cancel-order}
  Unfold the rule: with nothing paid it cancels the order directly.
\end{proof}

\begin{theorem}[Orders cancelled at once are closed]
  \label{thm:change-cancellation-cancelled-closed}
  \lean{Rules.OneOrder.Change.cancelled_at_once_closed}
  \uses{def:c15-cancel-order, def:c06-item-status, def:c07-order-status}
  Under every reading of the guide, an order cancelled at once has every item
  \texttt{CANCELLED} and is \texttt{CLOSED}.

  Source: \href{https://standards.atlassian.net/wiki/spaces/EASD/pages/574783569/Cancellation+of+Order+Items}{EASD guide, Cancellation of Order Items}.
\end{theorem}
\begin{proof}
  \uses{def:c15-cancel-order, def:c05-holds-value}
  Every service is cancelled and removed, and a removed service holds no value
  under any reading.
\end{proof}

\begin{theorem}[Penalty netted or charged separately (C-16)]
  \label{thm:c16-penalty-net}
  \lean{Rules.OneOrder.Change.net_same_either_way}
  \uses{def:c15-cancel-order, def:money, def:change-cancellation-due-amounts}
  A cancellation penalty is either deducted from the refund or charged
  separately, and the customer's net amount is the same either way.  For
  example, with 625.00 refundable and a penalty of 100.00, the customer is
  525.00 better off in both cases.  The model proves the two ways equal for
  changes and cancellations, under both readings of a penalty larger than the
  refundable total (C-85).

  Source: \href{https://standards.atlassian.net/wiki/spaces/EASD/pages/574783569/Cancellation+of+Order+Items}{EASD guide, Cancellation of Order Items}.

  Capabilities: ORDPEN.
\end{theorem}
\begin{proof}
  \uses{def:change-cancellation-due-amounts}
  Both ways, what the customer pays less what the airline returns is the price
  difference plus the penalty charged.
\end{proof}

\begin{definition}[Changing an order item (C-18)]
  \label{def:c18-change}
  \lean{Rules.OneOrder.Change.changeItem, Rules.OneOrder.Change.ItemChange, Rules.OneOrder.Change.originalAfterChange, Rules.OneOrder.Change.afterChange, Rules.OneOrder.Change.newItemOf}
  \uses{def:c01-structure, def:c06-item-status}
  Changing an order item cancels the original item's services, creates a new
  item, and charges or refunds the difference in price.  The model covers the
  change once it is paid: the services the change keeps move to the new item
  under the same references, their copies in the original item are marked
  \texttt{TRANSFERRED} with their delivery status unchanged, and the other
  services are cancelled and removed.  The money is priced by the price
  differential.

  Open point: the source shows the difference only in a worked example, and
  does not define its sign or how penalties enter it.

  Source: \href{https://standards.atlassian.net/wiki/spaces/EASD/pages/574619782/Modification+of+an+Order}{EASD guide, Modification of an Order}.

  Capabilities: ORDRSH.
\end{definition}

\begin{theorem}[Kept services keep their references]
  \label{thm:change-cancellation-kept-references}
  \lean{Rules.OneOrder.Change.change_keeps_services}
  \uses{def:c18-change}
  Every service a change keeps has a copy in the new item with the same
  reference, passenger, segment and delivery status.

  Source: \href{https://standards.atlassian.net/wiki/spaces/EASD/pages/574619782/Modification+of+an+Order}{EASD guide, Modification of an Order}.
\end{theorem}
\begin{proof}
  \uses{def:c18-change}
  The new item takes the kept services of the changed item and confirms them.
\end{proof}

\begin{theorem}[Originals of kept services are transferred]
  \label{thm:change-cancellation-originals-transferred}
  \lean{Rules.OneOrder.Change.change_marks_transferred}
  \uses{def:c18-change}
  Every service a change keeps stays in the changed item, marked
  \texttt{TRANSFERRED}, with its delivery status unchanged.

  Source: \href{https://standards.atlassian.net/wiki/spaces/EASD/pages/574619782/Modification+of+an+Order}{EASD guide, Modification of an Order}.
\end{theorem}
\begin{proof}
  \uses{def:c18-change}
  The changed item marks each service it keeps as transferred and changes
  nothing else about it.
\end{proof}

\begin{theorem}[Services not kept are cancelled]
  \label{thm:change-cancellation-others-cancelled}
  \lean{Rules.OneOrder.Change.change_cancels_others}
  \uses{def:c18-change}
  Every service of the changed item that the change does not keep is
  cancelled and removed.

  Source: \href{https://standards.atlassian.net/wiki/spaces/EASD/pages/574619782/Modification+of+an+Order}{EASD guide, Modification of an Order}.
\end{theorem}
\begin{proof}
  \uses{def:c18-change}
  The changed item marks each service it does not keep as cancelled and
  removed.
\end{proof}

\begin{theorem}[Changes keep item references distinct]
  \label{thm:change-cancellation-refs-distinct}
  \lean{Rules.OneOrder.Change.change_item_refs_unique}
  \uses{def:c18-change, thm:c03-unique-item-ids}
  If an order's item references are distinct before a change, they are
  distinct after it.  A change whose new item would reuse a reference is
  refused.

  Source: \href{https://standards.atlassian.net/wiki/spaces/EASD/pages/574619782/Modification+of+an+Order}{EASD guide, Modification of an Order}.
\end{theorem}
\begin{proof}
  \uses{def:c18-change}
  The changed item keeps its reference, and the new item's reference is not
  already in the order.
\end{proof}

\begin{theorem}[Transferred services keep the order balanced (C-19)]
  \label{thm:c19-transfer-balanced}
  \uses{def:c18-change, def:c06-item-status}
  \notready
  Services kept through a change move to the new item with the same service
  ID, and the originals are marked \texttt{TRANSFERRED}.  Delivered services
  move too, so that the order stays balanced.

  Open point: the source does not define what keeps an order balanced.

  Source: \href{https://standards.atlassian.net/wiki/spaces/EASD/pages/574849062/Order}{EASD guide, Order}.

  Capabilities: ORDRSH.
\end{theorem}

\begin{theorem}[Stored value stays in its order (C-20)]
  \label{thm:c20-value-store}
  \uses{def:c18-change, def:c05-holds-value}
  \notready
  Value stored as an order item can only be used in a later change to the same
  order.

  Open point: the source does not cover using part of the stored value, or
  when it expires.

  Source: \href{https://standards.atlassian.net/wiki/spaces/EASD/pages/574554245/Value+Store}{EASD guide, Value Store}.

  Capabilities: ORDREU.
\end{theorem}

% Second survey round: candidates written from the survey notes in docs/one-order-survey.md.

\begin{definition}[Price differential of a change (C-81)]
  \label{def:change-cancellation-price-differential}
  \lean{Rules.OneOrder.Change.priceDifferential, Rules.OneOrder.Change.ChangeQuote, Rules.OneOrder.Change.PriceDifferential, Rules.OneOrder.Change.DifferentialType, Rules.OneOrder.Change.statedDue, Rules.OneOrder.Change.typeOf}
  \uses{def:money, def:c18-change, def:c15-cancel-order, def:change-cancellation-due-amounts}
  Every change or cancellation offer carries a price differential: the old
  price, the new price (zero for a cancellation), the penalties, the
  forfeited amounts, and two non-negative amounts, one due to the airline
  and one due by the airline. A type code names the outcome: EvenExchange,
  AddCol, Refund, Residual (a value store), Reusable, AddColAndRefund or
  AddColAndResidual.

  Open point: The sign of the differential total is inconsistent. The
  Cancellation of Order Items page shows a refund with a positive total of
  625.00, while the worked examples show refunds with negative totals, such
  as minus 1,000.00 in \texttt{EXM\_ORD\_030A}. The business must fix one
  sign rule, or model only the two non-negative due amounts. The model states
  only the two due amounts until the business decides.

  Sources: \href{https://standards.atlassian.net/wiki/spaces/EASD/pages/574783569}{EASD guide, Cancellation of Order Items};
  \href{https://standards.atlassian.net/wiki/spaces/EASD/pages/842268673}{EASD guide, Customer requests a change to an Order containing a paid OrderItem};
  \href{https://standards.atlassian.net/wiki/spaces/EASD/pages/817594369}{EASD guide, Customer requests a cancellation to paid Order or OrderItems};
  \href{https://retailing.iata.org/armi/docs/ORDRSH/}{IATA capability ORDRSH};
  \href{https://standards.atlassian.net/wiki/spaces/EASD/pages/574554213/Price+Differentials}{EASD guide, Price Differentials}.

  Capabilities: ORDRSH, ORDCAN, ORDCA2, ORDPEN, ORDFOR, ORDREU.
\end{definition}

\begin{definition}[Amounts due to and by the airline]
  \label{def:change-cancellation-due-amounts}
  \lean{Rules.OneOrder.Change.dueAmounts, Rules.OneOrder.Change.DueAmounts, Rules.OneOrder.Change.ChangeQuote.difference, Rules.OneOrder.Change.ChangeQuote.refundable, Rules.OneOrder.Change.ExcessPenalty.charged, Rules.OneOrder.Change.PenaltyHandling, Rules.OneOrder.Change.Settlement}
  \uses{def:money, def:change-cancellation-penalty-exceeds-refund}
  The price difference of a change is the new price, less the old price, plus
  the forfeited amount.  With a netted penalty, the penalty charged is added
  to the difference, and the result is due to the airline when positive and
  by the airline when negative.  With a separate penalty, the penalty charged
  is due to the airline on its own, and the difference is settled on its own.
  The worked examples fit this formula; the guide does not state it.

  Source: \href{https://standards.atlassian.net/wiki/spaces/EASD/pages/574554213/Price+Differentials}{EASD guide, Price Differentials}.
\end{definition}

\begin{theorem}[Mixed currencies are rejected]
  \label{thm:change-cancellation-mixed-currency}
  \lean{Rules.OneOrder.Change.differential_rejects_mixed_currency}
  \uses{def:change-cancellation-price-differential}
  A change whose prices, penalty and forfeited amount are not all in one
  currency is rejected unpriced.

  Source: \href{https://standards.atlassian.net/wiki/spaces/EASD/pages/574619782/Modification+of+an+Order}{EASD guide, Modification of an Order}.
\end{theorem}
\begin{proof}
  \uses{def:change-cancellation-price-differential}
  The rule checks the currencies first.
\end{proof}

\begin{theorem}[Type code names a collection]
  \label{thm:change-cancellation-type-collects}
  \lean{Rules.OneOrder.Change.differential_collects_iff}
  \uses{def:change-cancellation-price-differential}
  A differential's type code is AddCol, AddColAndRefund or AddColAndResidual
  exactly when an amount is due to the airline.

  Source: \href{https://standards.atlassian.net/wiki/spaces/EASD/pages/574554213/Price+Differentials}{EASD guide, Price Differentials}.
\end{theorem}
\begin{proof}
  \uses{def:change-cancellation-price-differential}
  Check each branch of the type code against the due amounts.
\end{proof}

\begin{theorem}[Type code names a refund or a value store]
  \label{thm:change-cancellation-type-pays}
  \lean{Rules.OneOrder.Change.differential_pays_iff}
  \uses{def:change-cancellation-price-differential}
  A differential's type code is Refund, Residual, AddColAndRefund or
  AddColAndResidual exactly when an amount is due by the airline.

  Source: \href{https://standards.atlassian.net/wiki/spaces/EASD/pages/574554213/Price+Differentials}{EASD guide, Price Differentials}.
\end{theorem}
\begin{proof}
  \uses{def:change-cancellation-price-differential}
  Check each branch of the type code against the due amounts.
\end{proof}

\begin{theorem}[Value kept for later use is typed Reusable]
  \label{thm:change-cancellation-respend-reusable}
  \lean{Rules.OneOrder.Change.respend_reusable}
  \uses{def:change-cancellation-price-differential, def:change-cancellation-reusable-amount}
  When the seller asks to keep the value of cancelled items for later use, the
  differential is typed Reusable and both due amounts are zero.

  Source: \href{https://standards.atlassian.net/wiki/spaces/EASD/pages/817594921}{EASD guide, Customer requests a cancellation to paid Order or OrderItems resulting in a Reusable amount}.
\end{theorem}
\begin{proof}
  \uses{def:change-cancellation-price-differential}
  Unfold the rule: a request to keep the value states both amounts as zero.
\end{proof}

\begin{theorem}[Netted penalty gives one amount]
  \label{thm:change-cancellation-netted-one-amount}
  \lean{Rules.OneOrder.Change.netted_single_amount}
  \uses{def:change-cancellation-due-amounts}
  With a netted penalty, the due amounts net to the price difference plus the
  penalty charged, and at most one of them is above zero.

  Source: \href{https://standards.atlassian.net/wiki/spaces/EASD/pages/574783569/Cancellation+of+Order+Items}{EASD guide, Cancellation of Order Items}.
\end{theorem}
\begin{proof}
  \uses{def:change-cancellation-due-amounts}
  A netted result is split by its sign.
\end{proof}

\begin{theorem}[Separate penalty is added to the amount due]
  \label{thm:change-cancellation-separate-added}
  \lean{Rules.OneOrder.Change.separate_penalty_added}
  \uses{def:change-cancellation-due-amounts}
  With a separate penalty, the amount due to the airline is the penalty
  charged plus any price difference the customer pays, and the amount due by
  the airline is any price difference the airline returns.

  Source: \href{https://standards.atlassian.net/wiki/spaces/EASD/pages/574783569/Cancellation+of+Order+Items}{EASD guide, Cancellation of Order Items}.
\end{theorem}
\begin{proof}
  \uses{def:change-cancellation-due-amounts}
  Unfold the rule.
\end{proof}

\begin{theorem}[Due amounts net to the difference (C-82)]
  \label{thm:change-cancellation-due-net-identity}
  \lean{Rules.OneOrder.Change.differential_nets}
  \uses{def:change-cancellation-price-differential, thm:c16-penalty-net, def:c18-change, def:money}
  For every change or cancellation offer, the amount due to the airline
  minus the amount due by the airline equals the new price minus the old
  price, plus penalties, plus forfeited amounts. Netting of penalties or
  taxes changes how this net is split between the two due amounts and never
  changes the net. In \texttt{EXM\_ORD\_003C}, 900 minus 1,000 plus a 40
  penalty plus 120 forfeited gives 60 due to the airline.

  Case to check: The sample change record on the Modification of an Order
  page shows a Refund with old and new prices of zero and 100 due by the
  airline, in currency XSD against prices in USD; no penalty or forfeit
  explains it. A Reusable cancellation (\texttt{EXM\_ORD\_030E}) shows every
  amount as zero, so the identity holds only if the value kept on the
  cancelled item counts as the new price.

  Counterexample: a paid order at 1,000.00 cancelled with a request to keep
  the value for later use is typed Reusable with both amounts zero, where the
  formula gives 1,000.00 due by the airline (Lean example
  \texttt{respendCancellation}, after \texttt{EXM\_ORD\_030E}).  Under the
  reading that waives a penalty above the refundable value, a cancellation of
  1,000.00 with 20.00 forfeited and a 1,000.00 penalty has nothing due, where
  the formula gives 20.00 due to the airline (\texttt{penaltyAboveRefund}).

  Open point: The equation holds in all eleven worked examples that settle
  money, but the guide never states it, and C-18 leaves the sign and the
  penalty term open. The business must adopt this formula, or another one,
  as the definition of the amount owed on a change.

  Sources: \href{https://standards.atlassian.net/wiki/spaces/EASD/pages/842268745}{EASD guide, Customer requests a change of a paid order prior to travel resulting in additional collection due to penalty and forfeited amount};
  \href{https://standards.atlassian.net/wiki/spaces/EASD/pages/842268768}{EASD guide, Customer requests a change of a paid order prior to travel resulting in a refund};
  \href{https://standards.atlassian.net/wiki/spaces/EASD/pages/842301447}{EASD guide, Customer requests a change of a paid order prior to travel resulting in additional collection and a refund};
  \href{https://standards.atlassian.net/wiki/spaces/EASD/pages/817594543}{EASD guide, Customer requests a cancellation to paid Order or OrderItems resulting in Partial Refund with a Forfeited amount incurred};
  \href{https://standards.atlassian.net/wiki/spaces/EASD/pages/817594795}{EASD guide, Customer requests a cancellation to paid Order or OrderItems resulting in Partial Refund with a Forfeited amount incurred and requiring payment of a penalty};
  \href{https://standards.atlassian.net/wiki/spaces/EASD/pages/574619782}{EASD guide, Modification of an Order}.

  Capabilities: ORDRSH, ORDCAN, ORDPEN, ORDFOR.
\end{theorem}
\begin{proof}
  \uses{def:change-cancellation-price-differential}
  Not proven.  The counterexample above shows the statement is false for the
  rule as written.
\end{proof}

\begin{theorem}[Collecting the whole penalty keeps the formula]
  \label{thm:change-cancellation-net-of-collect}
  \lean{Rules.OneOrder.Change.due_amounts_net_of_collect}
  \uses{def:change-cancellation-due-amounts, def:change-cancellation-penalty-exceeds-refund}
  Under the reading that collects the whole penalty, the due amounts net to
  the new price, less the old price, plus the penalty, plus the forfeited
  amount, whether the penalty is netted or charged separately.

  Source: \href{https://standards.atlassian.net/wiki/spaces/EASD/pages/574783569/Cancellation+of+Order+Items}{EASD guide, Cancellation of Order Items}.
\end{theorem}
\begin{proof}
  \uses{def:change-cancellation-due-amounts}
  The penalty charged is the whole penalty, and either way the net is the
  price difference plus the penalty charged.
\end{proof}

\begin{theorem}[Order balance after servicing (C-83)]
  \label{thm:change-cancellation-order-balance}
  \lean{Rules.OneOrder.Change.change_keeps_balance}
  \uses{def:totals, def:payment, def:refund, thm:c19-transfer-balanced, def:c06-item-status, thm:change-cancellation-penalties-persist, thm:change-cancellation-due-net-identity, def:change-cancellation-ledger, def:c18-change}
  At every point in an order's life, payments minus refunds equal the prices
  of its live order items plus all recorded penalties plus all recorded
  forfeited amounts. Value store and reusable items count as live at their
  price. In \texttt{EXM\_ORD\_030B}, 1,000 paid minus 930 refunded equals a
  50 penalty plus 20 forfeited.

  Case to check: In \texttt{EXM\_ORD\_003C} after acceptance and before
  payment, the old item is still \texttt{ACTIVE} at 1,000 and the new item
  is \texttt{ACTIVE} at 900 with 1,000 paid, so the equation fails until the
  old item stops counting. In the Modification of an Order illustration,
  payments total 3,100 while the new item shows a price of 3,000 against a
  quoted new price of 3,100.

  Counterexample: a paid trip at 1,000.00 changed to a trip at 900.00 that
  keeps the outbound flight, with a 40.00 penalty and 120.00 forfeited, leaves
  1,060.00 paid against 1,900.00 of live items plus 160.00 recorded, because
  the original item still holds value.  The same change keeping nothing
  balances (Lean example \texttt{changeBalance}).

  Open point: The ORDPEN and ORDFOR guidance requires that sellers can
  reconcile all order values against payments and refunds but gives no
  equation, and C-19 leaves balance undefined. The business must choose the
  equation and say how an accepted but unpaid change enters it.

  Sources: \href{https://retailing.iata.org/armi/docs/ORDPEN/}{IATA capability ORDPEN};
  \href{https://retailing.iata.org/armi/docs/ORDFOR/}{IATA capability ORDFOR};
  \href{https://standards.atlassian.net/wiki/spaces/EASD/pages/842268745}{EASD guide, Customer requests a change of a paid order prior to travel resulting in additional collection due to penalty and forfeited amount};
  \href{https://standards.atlassian.net/wiki/spaces/EASD/pages/817594543}{EASD guide, Customer requests a cancellation to paid Order or OrderItems resulting in Partial Refund with a Forfeited amount incurred};
  \href{https://standards.atlassian.net/wiki/spaces/EASD/pages/574619782}{EASD guide, Modification of an Order};
  \href{https://retailing.iata.org/armi/docs/PAYREF/}{IATA capability PAYREF};
  \href{https://standards.atlassian.net/wiki/spaces/EASD/pages/842268768/Customer+requests+a+change+of+a+paid+order+prior+to+travel+resulting+in+a+refund}{EASD guide, Customer requests a change of a paid order prior to travel resulting in a refund};
  \href{https://standards.atlassian.net/wiki/spaces/EASD/pages/842301447/Customer+requests+a+change+of+a+paid+order+prior+to+travel+resulting+in+additional+collection+and+a+refund}{EASD guide, Customer requests a change of a paid order prior to travel resulting in additional collection and a refund};
  \href{https://standards.atlassian.net/wiki/spaces/EASD/pages/817594417/Customer+requests+a+cancellation+to+paid+Order+or+OrderItems+resulting+in+full+refund}{EASD guide, Customer requests a cancellation to paid Order or OrderItems resulting in full refund}.

  Capabilities: ORDPEN, ORDFOR, ORDRSH, PAYREF, PAYSUM.
\end{theorem}
\begin{proof}
  \uses{def:change-cancellation-ledger, def:c18-change}
  Not proven.  The counterexample above shows the statement is false for the
  rule as written.
\end{proof}

\begin{definition}[Money recorded on an order]
  \label{def:change-cancellation-ledger}
  \lean{Rules.OneOrder.Change.settle, Rules.OneOrder.Change.Ledger, Rules.OneOrder.Change.liveValue, Rules.OneOrder.Change.Ledger.Balanced}
  \uses{def:money, def:change-cancellation-price-differential, def:c05-holds-value}
  An order records everything the customer paid, everything the airline
  returned, and every penalty and forfeited amount, at order level.  Settling
  a change or cancellation offer adds its amount due to the airline to the
  payments, its amount due by the airline to the refunds, and its penalty and
  forfeited amount to those recorded.  The live items are the order's active
  items.

  Source: \href{https://retailing.iata.org/armi/docs/ORDPEN/}{IATA capability ORDPEN}.
\end{definition}

\begin{theorem}[Settling adds to payments and refunds]
  \label{thm:change-cancellation-settle-payments}
  \lean{Rules.OneOrder.Change.settle_adds_payments}
  \uses{def:change-cancellation-ledger}
  Settling an offer adds the amount due to the airline to the payments and the
  amount due by the airline to the refunds.

  Source: \href{https://retailing.iata.org/armi/docs/ORDPEN/}{IATA capability ORDPEN}.
\end{theorem}
\begin{proof}
  \uses{def:change-cancellation-ledger}
  Unfold the rule.
\end{proof}

\begin{theorem}[Settling records penalties and forfeits]
  \label{thm:change-cancellation-settle-records}
  \lean{Rules.OneOrder.Change.settle_records_penalty}
  \uses{def:change-cancellation-ledger}
  Settling an offer adds its penalty and forfeited amount to those recorded on
  the order, and lowers no amount recorded before.

  Source: \href{https://retailing.iata.org/armi/docs/ORDPEN/}{IATA capability ORDPEN}.
\end{theorem}
\begin{proof}
  \uses{def:change-cancellation-ledger}
  Settling only adds.
\end{proof}

\begin{definition}[Change accepted before payment (C-84)]
  \label{def:change-cancellation-deferred-payment-change}
  \uses{def:c11-payment-time-limit, def:c08-delivery-status, def:c05-holds-value, def:c18-change}
  \notready
  A seller may accept a change offer and pay later. Until payment, the new
  item's services are \texttt{CONFIRMED} with a payment time limit, and the
  old item stays \texttt{ACTIVE} with its replaced services
  \texttt{CANCELLED} and in \texttt{REMOVED HOLD VALUE}. Payment moves the
  new services to \texttt{READY TO PROCEED} and the old item to
  \texttt{CANCELLED}.

  Case to check: In \texttt{EXM\_ORD\_003C} the change is accepted late on
  18 May 2023 with payment due by 16:00 UTC on 19 May. If the 60.00 is never
  paid, neither item gives the customer a flight.

  Open point: The source does not say what happens if the payment time limit
  passes. The business must decide whether the old services return to
  \texttt{READY TO PROCEED}, or the customer is left with an unpaid new item
  and cancelled old services.

  Sources: \href{https://standards.atlassian.net/wiki/spaces/EASD/pages/842268673}{EASD guide, Customer requests a change to an Order containing a paid OrderItem};
  \href{https://standards.atlassian.net/wiki/spaces/EASD/pages/842268745}{EASD guide, Customer requests a change of a paid order prior to travel resulting in additional collection due to penalty and forfeited amount};
  \href{https://standards.atlassian.net/wiki/spaces/EASD/pages/842301447}{EASD guide, Customer requests a change of a paid order prior to travel resulting in additional collection and a refund};
  \href{https://standards.atlassian.net/wiki/spaces/EASD/pages/842301447/Customer+requests+a+change+of+a+paid+order+prior+to+travel+resulting+in+additional+collection+and+a+refund}{EASD guide, Customer requests a change of a paid order prior to travel resulting in additional collection and a refund};
  \href{https://standards.atlassian.net/wiki/spaces/EASD/pages/842268768/Customer+requests+a+change+of+a+paid+order+prior+to+travel+resulting+in+a+refund}{EASD guide, Customer requests a change of a paid order prior to travel resulting in a refund}.

  Capabilities: ORDRSH, ORDPEN, ORDFOR, PAYREF, PAYORD.
\end{definition}

\begin{definition}[Penalty larger than the refund (C-85)]
  \label{def:change-cancellation-penalty-exceeds-refund}
  \lean{Rules.OneOrder.Change.ExcessPenalty}
  \uses{def:money, def:c15-cancel-order}
  With a netted penalty, the amount due by the airline is the refundable
  value of the cancelled items minus the penalty. The source has no case
  where the penalty is larger than the refundable value.

  Case to check: Take \texttt{EXM\_ORD\_030B} with a cancellation fee of
  1,000 in place of 50. The refundable value is 980, so the netted result is
  20 due to the airline on an offer typed Refund.

  Open point: The business must decide whether a penalty is capped at the
  refundable value, so that cancelling never leaves the customer owing
  money, or whether the excess is collected. The model takes the answer as a
  parameter with two readings: collect the excess, or waive it.

  Sources: \href{https://standards.atlassian.net/wiki/spaces/EASD/pages/574783569}{EASD guide, Cancellation of Order Items};
  \href{https://standards.atlassian.net/wiki/spaces/EASD/pages/817594543}{EASD guide, Customer requests a cancellation to paid Order or OrderItems resulting in Partial Refund with a Forfeited amount incurred}.

  Capabilities: ORDCAN, ORDCA2, ORDPEN.
\end{definition}

\begin{theorem}[Cancelling leaves the customer owing nothing]
  \label{thm:change-cancellation-cancel-not-owing}
  \lean{Rules.OneOrder.Change.cancel_not_owing}
  \uses{def:change-cancellation-due-amounts, def:change-cancellation-penalty-exceeds-refund}
  On a cancellation whose forfeited amount is part of the old price, the
  amount due to the airline is no larger than the amount due by the airline,
  under both readings of a penalty above the refundable value.

  Counterexample: a paid order at 1,000.00 with 20.00 forfeited, cancelled
  with a netted penalty of 1,000.00, leaves 20.00 due to the airline under the
  reading that collects the whole penalty (Lean example
  \texttt{penaltyAboveRefund}).

  Open point: under the other reading the guarantee holds, but the due
  amounts no longer net to the formula of C-82.  The business must decide
  whether a penalty is capped at the refundable value or the excess is
  collected.

  Source: \href{https://standards.atlassian.net/wiki/spaces/EASD/pages/574783569/Cancellation+of+Order+Items}{EASD guide, Cancellation of Order Items}.
\end{theorem}
\begin{proof}
  \uses{def:change-cancellation-due-amounts}
  Not proven.  The counterexample above shows the statement is false for the
  rule as written.
\end{proof}

\begin{theorem}[Waiving the excess penalty leaves nothing owing]
  \label{thm:change-cancellation-not-owing-of-waive}
  \lean{Rules.OneOrder.Change.cancel_not_owing_of_waive}
  \uses{def:change-cancellation-due-amounts, def:change-cancellation-penalty-exceeds-refund}
  Under the reading that waives a penalty above the refundable value, a
  cancellation never leaves the customer owing the airline money.

  Source: \href{https://standards.atlassian.net/wiki/spaces/EASD/pages/574783569/Cancellation+of+Order+Items}{EASD guide, Cancellation of Order Items}.
\end{theorem}
\begin{proof}
  \uses{def:change-cancellation-due-amounts}
  The penalty charged is at most the refundable value, and a cancellation
  returns the refundable value.
\end{proof}

\begin{theorem}[Full cancellation keeps no residual (C-86)]
  \label{thm:change-cancellation-full-cancel-no-residual}
  \uses{def:c15-cancel-order, thm:c20-value-store, def:c07-order-status}
  \notready
  Cancelling a whole order cancels every service and settles any refundable
  amount. The customer cannot keep residual value in the order for a later
  reshop.

  Case to check: If a full cancellation must close the order, 030C conflicts
  with C-07, because its value store service still holds value and keeps the
  order \texttt{OPENED}.

  Open point: The ORDCA2 guidance states this, but worked example
  \texttt{EXM\_ORD\_030C} cancels a whole order through the full
  cancellation request and keeps the full 1,000 as a Value Store order item,
  leaving the order \texttt{OPENED}. The business must decide whether a full
  cancellation may end in a value store or a reusable amount.

  Sources: \href{https://retailing.iata.org/armi/docs/ORDCA2/}{IATA capability ORDCA2};
  \href{https://standards.atlassian.net/wiki/spaces/EASD/pages/817594669}{EASD guide, Customer requests a cancellation to paid Order or OrderItems resulting in Value Store}.

  Capabilities: ORDCA2, ORDREU.
\end{theorem}

\begin{definition}[Paid or partly paid orders need an offer (C-87)]
  \label{def:change-cancellation-partly-paid-cancel}
  \lean{Rules.OneOrder.Change.paid_needs_offer}
  \uses{def:c15-cancel-order, thm:refunds-le-payments, def:payment}
  An order with any payment, full or partial, is cancelled through a
  cancellation offer. Only an order with no payment at all may be cancelled
  directly.

  Case to check: Take an order priced 1,000 with 300 paid. The worked
  examples compute refundable value from the old price, which would show
  1,000 due by the airline against 300 received and break the rule that
  refunds never exceed payments.

  Open point: The EASD page speaks only of paid and unpaid orders; the
  partial case comes from ORDCA2. The business must decide which route
  applies when some items are paid and others are not, and whether the
  refundable amount is capped at the amount received.

  Sources: \href{https://retailing.iata.org/armi/docs/ORDCA2/}{IATA capability ORDCA2};
  \href{https://standards.atlassian.net/wiki/spaces/EASD/pages/574783569}{EASD guide, Cancellation of Order Items}.

  Capabilities: ORDCA2, ORDCAN.
\end{definition}

\begin{theorem}[Cancellation offers stay within payments]
  \label{thm:change-cancellation-offer-within-payments}
  \lean{Rules.OneOrder.Change.offer_within_payments}
  \uses{def:c15-cancel-order, def:change-cancellation-partly-paid-cancel, def:change-cancellation-price-differential, def:totals}
  A cancellation offer never makes the amount due by the airline larger than
  the payments received for the order.

  Counterexample: an order at 1,000.00 with 300.00 received is cancelled
  through an offer priced from the order's price, with 1,000.00 due by the
  airline (Lean example \texttt{partlyPaidCancellation}).

  Open point: the business must decide whether the refundable amount is
  capped at the payments received, and which route applies when some items
  are paid and others are not.

  Source: \href{https://retailing.iata.org/armi/docs/ORDCA2/}{IATA capability ORDCA2}.
\end{theorem}
\begin{proof}
  \uses{def:c15-cancel-order, def:change-cancellation-price-differential}
  Not proven.  The counterexample above shows the statement is false for the
  rule as written.
\end{proof}

\begin{definition}[Removing one passenger from an item (C-88)]
  \label{def:change-cancellation-remove-passenger}
  \uses{thm:c04-item-unsplit, thm:c19-transfer-balanced, def:c18-change, def:passenger}
  \notready
  When one order item covers several passengers, removing a passenger
  deletes the whole item and lists the other passengers' services to retain.
  The retained services move to a new item, which is repriced.

  Case to check: Two passengers share one item priced 1,000 under a
  companion fare. Removing one reprices the other at a 900 single fare, so
  the leaving passenger is refunded 100 where an even split gives 500.

  Open point: The source does not say whether the retained services keep
  their share of the original price or are repriced at current fares, or how
  a penalty is assigned when one passenger leaves. The same page also
  classes cancelling services inside an item as a change of order, yet uses
  that method for passenger removal under cancellation.

  Sources: \href{https://standards.atlassian.net/wiki/spaces/EASD/pages/574783569}{EASD guide, Cancellation of Order Items};
  \href{https://retailing.iata.org/armi/docs/ORDCAN/}{IATA capability ORDCAN}.

  Capabilities: ORDCAN, ORDRSH.
\end{definition}

\begin{definition}[Reusable amount on cancellation (C-89)]
  \label{def:change-cancellation-reusable-amount}
  \uses{def:c05-holds-value, def:c06-item-status, def:c07-order-status, thm:c20-value-store}
  \notready
  A cancellation that asks to respend (RespendInd true) triggers no refund.
  The item stays \texttt{ACTIVE} at its full price, its services become
  \texttt{CANCELLED} with delivery status \texttt{REMOVED HOLD VALUE}, the
  differential is typed Reusable with zero amounts, and the order stays
  \texttt{OPENED}. The exact amount is computed later, at a refund request
  or a reshop.

  Case to check: In \texttt{EXM\_ORD\_030E} the customer paid 1,000 and
  holds an \texttt{ACTIVE} item at 1,000. If the airline recomputes the
  value lower at reuse, nothing in the source protects the 1,000.

  Open point: The source sets no floor, ceiling or time limit on the amount
  computed later. The business must decide whether the customer is owed at
  least what a refund would have paid on the day of cancellation, and when
  the held value lapses.

  Sources: \href{https://retailing.iata.org/armi/docs/ORDREU/}{IATA capability ORDREU};
  \href{https://standards.atlassian.net/wiki/spaces/EASD/pages/817594921}{EASD guide, Customer requests a cancellation to paid Order or OrderItems resulting in a Reusable amount};
  \href{https://standards.atlassian.net/wiki/spaces/EASD/pages/574783569}{EASD guide, Cancellation of Order Items}.

  Capabilities: ORDREU, ORDCAN.
\end{definition}

\begin{definition}[Settling money the airline owes (C-90)]
  \label{def:change-cancellation-owed-settlement}
  \uses{def:refund, def:c17-refund-methods, thm:c20-value-store, def:change-cancellation-reusable-amount, def:change-cancellation-price-differential}
  \notready
  Money the airline owes after a change or cancellation is settled as a
  refund (by default to the original form of payment), a Value Store order
  item, a reusable amount held on the cancelled item, or value on an open
  accountable document. Non-refundable parts, such as a non-refundable tax,
  are recorded as forfeited.

  Case to check: In \texttt{EXM\_ORD\_003D} and 003E the refund offer and
  the value store offer are both worth 150. Nothing in the source stops a
  150 refund offer alongside a 120 value store offer.

  Open point: The sources differ on who picks the form. The change page has
  the airline choose between refund and value store, the cancellation page
  lets the seller request respend, and \texttt{EXM\_ORD\_003D} and 003E show
  the airline offering both while the seller accepts one. The business must
  decide who chooses, and whether a value store must equal the refund it
  replaces.

  Sources: \href{https://standards.atlassian.net/wiki/spaces/EASD/pages/574554168}{EASD guide, Seller Initiated Change to an Order};
  \href{https://standards.atlassian.net/wiki/spaces/EASD/pages/574783569}{EASD guide, Cancellation of Order Items};
  \href{https://standards.atlassian.net/wiki/spaces/EASD/pages/574554245}{EASD guide, Value Store};
  \href{https://retailing.iata.org/armi/docs/ORDCAN/}{IATA capability ORDCAN};
  \href{https://standards.atlassian.net/wiki/spaces/EASD/pages/842268768}{EASD guide, Customer requests a change of a paid order prior to travel resulting in a refund};
  \href{https://standards.atlassian.net/wiki/spaces/EASD/pages/842268791}{EASD guide, Customer requests a change of a paid order prior to travel resulting in a value store}.

  Capabilities: ORDCAN, ORDRSH, ORDREU, ORDFOR.
\end{definition}

\begin{definition}[Value store time windows (C-91)]
  \label{def:change-cancellation-value-store-windows}
  \uses{thm:c20-value-store, def:instant, def:c05-holds-value, def:c07-order-status}
  \notready
  A Value Store order item carries its own cancel and change restrictions
  with start and end dates. In \texttt{EXM\_ORD\_030C} the 1,000 value store
  may be cancelled at no fee for one year and used in a change for two
  years, and cancelling is not allowed after the first year.

  Case to check: A customer who holds the 030C value store for 25 months can
  neither cancel nor use it, yet the order stays \texttt{OPENED} with a
  service that still holds value under C-05.

  Open point: The source does not say what the customer receives on
  cancelling a value store, or what happens to unused value when the change
  window closes: refund, forfeit or a voucher. C-20 leaves expiry open;
  these windows supply the dates but not the outcome.

  Sources: \href{https://standards.atlassian.net/wiki/spaces/EASD/pages/817594669}{EASD guide, Customer requests a cancellation to paid Order or OrderItems resulting in Value Store};
  \href{https://standards.atlassian.net/wiki/spaces/EASD/pages/842268791}{EASD guide, Customer requests a change of a paid order prior to travel resulting in a value store};
  \href{https://standards.atlassian.net/wiki/spaces/EASD/pages/574554245}{EASD guide, Value Store}.

  Capabilities: ORDCA2, ORDCAN.
\end{definition}

\begin{theorem}[Stored value is used once (C-92)]
  \label{thm:change-cancellation-stored-value-used-once}
  \uses{thm:c20-value-store, def:c05-holds-value, def:change-cancellation-owed-settlement}
  \notready
  The total drawn from any stored value, whether a Value Store item, a
  reusable item or an open accountable document, never exceeds its value
  across all orders. Value moved to a new order leaves the order it came
  from.

  Case to check: A seller reallocates 500 from an open ticket to a new order
  through OrderCreate, then quotes a change on the original order naming the
  same ticket. Nothing in the source marks the 500 as spent.

  Open point: The Value Store page lets value on an open document be
  reallocated to a new order or within the changed order, recalculated on
  reuse, but states no rule against using it twice. The business must decide
  how stored value is tracked and when it is used up, including the partial
  use that C-20 leaves open.

  Source: \href{https://standards.atlassian.net/wiki/spaces/EASD/pages/574554245}{EASD guide, Value Store}.

  Capabilities: ORDREU.
\end{theorem}

\begin{theorem}[Transferred originals hold no value (C-93)]
  \label{thm:change-cancellation-transferred-original-no-value}
  \lean{Rules.OneOrder.Change.change_original_cancelled}
  \uses{def:c05-holds-value, def:c06-item-status, thm:c19-transfer-balanced, def:c08-delivery-status, def:c18-change}
  A service transferred to a new item holds value only in the new item. The
  original copy, with status \texttt{TRANSFERRED}, holds no value, so its
  old item can be \texttt{CANCELLED}.

  Case to check: In \texttt{EXM\_ORD\_003B} after the change, ORDITM01 is
  \texttt{CANCELLED} while six of its services show \texttt{TRANSFERRED}
  with \texttt{READY TO PROCEED}. The same six service IDs appear in
  \texttt{ORDITM02\_NEW} as \texttt{READY TO PROCEED}, so their value is
  counted twice.

  Counterexample: a paid return trip changed to keep the outbound flight
  leaves the original item with the outbound service \texttt{TRANSFERRED} and
  \texttt{READY TO PROCEED}.  Under both readings the original item is
  \texttt{ACTIVE}, and the outbound service holds value in both items (Lean
  example \texttt{changeKeepingOutbound}).

  Open point: In the worked examples the original copy keeps delivery status
  \texttt{READY TO PROCEED} while its item is \texttt{CANCELLED}. C-05
  judges value by delivery status alone, under which the cancelled item
  still holds value and meets the \texttt{ACTIVE} condition of C-06. The
  business must decide whether the \texttt{TRANSFERRED} status overrides the
  delivery status.

  Sources: \href{https://standards.atlassian.net/wiki/spaces/EASD/pages/841187358}{EASD guide, Customer requests a change of a paid order prior to travel resulting in additional collection};
  \href{https://retailing.iata.org/armi/docs/ORDRSH/}{IATA capability ORDRSH};
  \href{https://standards.atlassian.net/wiki/spaces/EASD/pages/574849062/Order}{EASD guide, Order};
  \href{https://standards.atlassian.net/wiki/spaces/OA/pages/844202130}{Order Accounting, Customer requests a change of a paid order prior to travel resulting in additional collection and a refund};
  \href{https://standards.atlassian.net/wiki/spaces/EASD/pages/842268712}{EASD guide, Customer requests a change of a paid order prior to travel resulting in additional collection and a value store}.

  Capabilities: ORDRSH.
\end{theorem}
\begin{proof}
  \uses{def:c18-change, def:c06-item-status}
  Not proven.  The counterexample above shows the statement is false for the
  rule as written.
\end{proof}

\begin{theorem}[Refund waits for separate penalty (C-94)]
  \label{thm:change-cancellation-refund-after-penalty}
  \uses{thm:c16-penalty-net, def:refund, def:payment, thm:c20-value-store}
  When a penalty is charged separately, the airline releases the refund, or
  makes the value store usable, only after the penalty is paid. A netted
  penalty leaves nothing to wait for.

  Case to check: C-16 treats netted and separate penalties as equivalent. In
  \texttt{EXM\_ORD\_003G} the separate penalty means the customer pays 60 in
  cash and receives a 160 value store; netted, the customer would pay
  nothing and receive 100 stored. The net is equal, but 60 of cash becomes
  restricted stored value.

  Sources: \href{https://standards.atlassian.net/wiki/spaces/EASD/pages/817594369}{EASD guide, Customer requests a cancellation to paid Order or OrderItems};
  \href{https://standards.atlassian.net/wiki/spaces/EASD/pages/842301447}{EASD guide, Customer requests a change of a paid order prior to travel resulting in additional collection and a refund};
  \href{https://standards.atlassian.net/wiki/spaces/EASD/pages/842268712}{EASD guide, Customer requests a change of a paid order prior to travel resulting in additional collection and a value store};
  \href{https://standards.atlassian.net/wiki/spaces/OA/pages/844169238}{Order Accounting, Customer requests a change of a paid order prior to travel resulting in additional collection and a value store}.

  Capabilities: ORDPEN, ORDCAN.
\end{theorem}

\begin{theorem}[Penalties and forfeits persist (C-95)]
  \label{thm:change-cancellation-penalties-persist}
  \uses{def:order, def:c18-change, def:money}
  From release 24.1, penalties and forfeited amounts are recorded at order
  level with links to order items. Once recorded, each keeps its amount and
  stays on the order when its order item is later replaced, cancelled or
  dropped from the order view.

  Case to check: In the change worked examples the change penalty points at
  the new item \texttt{ORDITM02\_NEW}, and the next retrieval no longer
  shows the cancelled item. A model that stores penalties on items would
  lose them at the next change.

  Sources: \href{https://retailing.iata.org/armi/docs/ORDPEN/}{IATA capability ORDPEN};
  \href{https://retailing.iata.org/armi/docs/ORDFOR/}{IATA capability ORDFOR};
  \href{https://standards.atlassian.net/wiki/spaces/EASD/pages/842268673}{EASD guide, Customer requests a change to an Order containing a paid OrderItem}.

  Capabilities: ORDPEN, ORDFOR.
\end{theorem}

\begin{definition}[Name change through an offer (C-96)]
  \label{def:change-cancellation-name-change}
  \uses{def:passenger, thm:delivered-not-cancellable, def:c08-delivery-status, def:c06-item-status}
  \notready
  A name change or correction is made by accepting a name change offer, free
  or paid. It adds a name change order item priced at the fee, whose service
  is \texttt{DELIVERED} at once; it overwrites all name elements of the same
  passenger and leaves the flight items and their services unchanged.
  Birthdate, gender, birthplace and additional names cannot be changed this
  way.

  Case to check: In the paid name change example the fee item ORDITM02 is
  \texttt{CLOSED} at 100 and the flight item ORDITM01 is unchanged. A full
  cancellation the next day would refund the flight item under its own rules
  and keep the 100.

  Open point: The source does not separate correcting a name from handing
  the booking to another person. It also does not say whether a paid name
  change fee is refunded if the order is later cancelled; because its
  service is \texttt{DELIVERED}, the existing rule forbids cancelling it.

  Sources: \href{https://retailing.iata.org/armi/docs/ORDNAM/}{IATA capability ORDNAM};
  \href{https://standards.atlassian.net/wiki/spaces/EASD/pages/672825346}{EASD guide, Customer Changes Their Name};
  \href{https://standards.atlassian.net/wiki/spaces/EASD/pages/672792581}{EASD guide, Paid name changes}.

  Capabilities: ORDNAM.
\end{definition}

\begin{definition}[Changes made without an offer (C-97)]
  \label{def:change-cancellation-no-offer-changes}
  \uses{def:c15-cancel-order, thm:c21-version, def:passenger, def:change-cancellation-partly-paid-cancel}
  \notready
  A change goes through an offer when it replaces, adds or cancels items of
  an order with any payment, removes a passenger, or changes a name. Other
  passenger details, sharing preferences, cancelling an unpaid order and
  accepting a follow-up action are made directly; they keep the passenger ID
  and have no payment step.

  Open point: The direct path has no price, yet the sample change record on
  the Modification of an Order page pairs an Accept follow-up action with
  100 due by the airline. The business must decide whether a direct change
  can create an amount owed and, if so, how it is priced.

  Sources: \href{https://standards.atlassian.net/wiki/spaces/EASD/pages/574554168}{EASD guide, Seller Initiated Change to an Order};
  \href{https://retailing.iata.org/armi/docs/ORDPAX/}{IATA capability ORDPAX};
  \href{https://standards.atlassian.net/wiki/spaces/EASD/pages/574619782}{EASD guide, Modification of an Order}.

  Capabilities: ORDPAX, ORDRSH, ORDNAM, ORDCAN, ORDCA2.
\end{definition}

\begin{theorem}[Full cancellation equals item cancellations (C-98)]
  \label{thm:change-cancellation-full-equals-items}
  \uses{def:c15-cancel-order, thm:c16-penalty-net, thm:change-cancellation-full-cancel-no-residual}
  \notready
  Cancelling a whole order produces the same refundable amounts and
  penalties as cancelling each of its items under the same rules.

  Case to check: An order has two items and a cancellation fee of 50 per
  transaction. Cancelling the order costs 50; cancelling the items one at a
  time costs 100.

  Open point: ORDCA2 says the calculation mechanisms are identical but does
  not say the totals must match. The business must decide whether
  order-level fees apply once per transaction or once per item, and whether
  item-by-item cancellation may keep reusable value that a full cancellation
  may not.

  Source: \href{https://retailing.iata.org/armi/docs/ORDCA2/}{IATA capability ORDCA2}.

  Capabilities: ORDCA2, ORDCAN.
\end{theorem}

\begin{theorem}[Stored value belongs to payer (C-99)]
  \label{thm:change-cancellation-stored-value-payer}
  \uses{def:payment, def:c13-several-methods, def:c17-refund-methods, thm:c20-value-store}
  \notready
  A value store is held by the airline for the payer of the transactions
  that created it.

  Open point: The source does not say whose value it is when several payers
  or payment methods paid the original item, or when the payer is not the
  passenger. The business must decide who may use, move or reclaim it.

  Source: \href{https://standards.atlassian.net/wiki/spaces/EASD/pages/574554245}{EASD guide, Value Store}.

  Capabilities: ORDREU.
\end{theorem}

% Candidate rules for airline-initiated change.  Written by hand; see docs/one-order-coverage.md.

\section{Airline-initiated change}

% Second survey round: candidates written from the survey notes in docs/one-order-survey.md.

\begin{theorem}[Penalty when declining an airline change (C-100)]
  \label{thm:airline-change-cancel-after-involuntary}
  \uses{thm:c16-penalty-net, def:c15-cancel-order, def:airline-change-follow-up-action, def:airline-change-involuntary-change}
  \notready
  A customer who declines an airline-initiated change through the Cancel
  follow-up action pays no cancellation penalty.

  Case to check: In the itinerary change worked example the new order item
  carries a 50 EUR cancellation fee. Applying it to a Cancel follow-up
  charges the customer for the airline's own change.

  Open point: The sources offer Cancel as a response to an airline change
  but never say whether penalties apply. The cancellation page covers only
  customer-initiated cancellation and leaves penalties to the airline.

  Sources: \href{https://standards.atlassian.net/wiki/spaces/EASD/pages/574849094}{EASD guide, Seller Follow-up Actions};
  \href{https://standards.atlassian.net/wiki/spaces/EASD/pages/672825400}{EASD guide, Inform of Impacted Dependent Services};
  \href{https://standards.atlassian.net/wiki/spaces/EASD/pages/574783569}{EASD guide, Cancellation of Order Items};
  \href{https://standards.atlassian.net/wiki/spaces/EASD/pages/815202352}{EASD guide, Airline notifies Seller of changes to an Order}.
\end{theorem}

\begin{theorem}[Refund for a cancelled flight (C-101)]
  \label{thm:airline-change-refund-undeliverable}
  \uses{def:c05-holds-value, def:c08-delivery-status, thm:c04-item-unsplit, def:refund, def:ancillaries-flight-change-dependents}
  \notready
  When the airline cancels a flight and does not re-accommodate, the flight
  service and its dependent services become \texttt{CANCELLED} with delivery
  status \texttt{UNABLE TO DELIVER}. The customer is then owed a refund of
  the value of those services.

  Case to check: A two-segment item costs 1000 EUR and the airline cancels
  the second segment. The item stays \texttt{ACTIVE} because the first
  segment still holds value, and no rule returns any part of the 1000 EUR.

  Open point: The worked example offers no follow-up action and no refund,
  and no source says whether \texttt{UNABLE TO DELIVER} caused by the
  airline entitles the customer to money back. Because price is set per
  order item, the business must also decide how much of an item's price
  belongs to one cancelled segment.

  Sources: \href{https://standards.atlassian.net/wiki/spaces/EASD/pages/815202352}{EASD guide, Airline notifies Seller of changes to an Order};
  \href{https://standards.atlassian.net/wiki/spaces/EASD/pages/574849062}{EASD guide, Order}.
\end{theorem}

\begin{definition}[Consequence of a missed time limit (C-102)]
  \label{def:airline-change-inaction-consequence}
  \uses{def:airline-change-follow-up-action, def:instant, def:c15-cancel-order, def:refund}
  \notready
  If a follow-up action carries a time limit and the seller has not acted by
  then, the airline applies the stated consequence. Auto-acceptance keeps
  the changed order and Auto-cancellation cancels it. The time limit and the
  consequence are both optional.

  Case to check: The concept page pairs Accept with Auto-cancellation. A
  seller who misses the deadline then costs the customer the whole trip
  after a change the airline made.

  Open point: The business must decide what applies when either the time
  limit or the consequence is missing (the Cancel example has neither), and
  what Auto-cancellation cancels and refunds. It must also decide whether
  the time limit has to fall before the affected flight departs.

  Sources: \href{https://standards.atlassian.net/wiki/spaces/EASD/pages/574849094}{EASD guide, Seller Follow-up Actions};
  \href{https://standards.atlassian.net/wiki/spaces/EASD/pages/815202352}{EASD guide, Airline notifies Seller of changes to an Order};
  \href{https://standards.atlassian.net/wiki/spaces/EASD/pages/672825400}{EASD guide, Inform of Impacted Dependent Services}.
\end{definition}

\begin{theorem}[Change applied before seller response (C-103)]
  \label{thm:airline-change-applied-before-acceptance}
  \uses{def:c06-item-status, thm:c21-version, thm:c19-transfer-balanced, def:airline-change-follow-up-action}
  \notready
  An airline-initiated change takes effect in the order when it is notified.
  The version rises, replaced services are cancelled or transferred, and the
  new order item is \texttt{ACTIVE} while the seller's Accept is still
  pending.

  Case to check: In the itinerary example the original item is already
  \texttt{CANCELLED} at version 2 while Accept is pending. A customer who
  rejects the change cannot get the original flights back, so Accept offers
  no real choice.

  Open point: This is shown only in a worked example. The business must
  decide whether a change awaiting acceptance is provisional, and whether
  declining it restores the original item or only lets the customer cancel
  or reshop from the changed order.

  Sources: \href{https://standards.atlassian.net/wiki/spaces/EASD/pages/815202352}{EASD guide, Airline notifies Seller of changes to an Order};
  \href{https://standards.atlassian.net/wiki/spaces/EASD/pages/574554168}{EASD guide, Seller Initiated Change to an Order}.
\end{theorem}

\begin{theorem}[Paid status after re-accommodation (C-104)]
  \label{thm:airline-change-reaccommodated-paid-status}
  \uses{def:c08-delivery-status, def:c11-payment-time-limit, thm:c19-transfer-balanced, thm:airline-change-applied-before-acceptance}
  \notready
  A service the airline creates to replace a paid service, with no price
  difference owed, starts with delivery status \texttt{READY TO PROCEED}.

  Case to check: If \texttt{CONFIRMED} is read as unpaid, a payment time
  limit could apply to a re-accommodated trip the customer already paid for
  in full.

  Open point: The itinerary worked example gives such services delivery
  status \texttt{CONFIRMED}, which the Order concept reserves for services
  still awaiting payment, although the exchange is even. The business must
  decide whether paid status carries across an airline-initiated change.

  Sources: \href{https://standards.atlassian.net/wiki/spaces/EASD/pages/815202352}{EASD guide, Airline notifies Seller of changes to an Order};
  \href{https://standards.atlassian.net/wiki/spaces/EASD/pages/574849062}{EASD guide, Order};
  \href{https://standards.atlassian.net/wiki/spaces/EASD/pages/672825400}{EASD guide, Inform of Impacted Dependent Services}.
\end{theorem}

\begin{theorem}[Unacknowledged change notice (C-105)]
  \label{thm:airline-change-lost-notice}
  \uses{def:airline-change-inaction-consequence, def:airline-change-notice-chain}
  \notready
  A follow-up consequence takes effect only after the seller has
  acknowledged the notice that carried it.

  Case to check: A notice carrying Accept with Auto-cancellation is lost.
  The seller sees the pending action only if it retrieves the order, and the
  order may be cancelled before it does.

  Open point: The guide names an unreceived or unacknowledged notice as the
  failure case but attaches no rule to it: no retry, no fallback, and no
  effect on the time limit. The business must decide whether the time limit
  runs from sending or from acknowledgement.

  Sources: \href{https://standards.atlassian.net/wiki/spaces/EASD/pages/574455903}{EASD guide, Notification of Airline Initiated Change to an Order};
  \href{https://standards.atlassian.net/wiki/spaces/EASD/pages/574849094}{EASD guide, Seller Follow-up Actions}.
\end{theorem}

\begin{definition}[Changes that require a customer choice (C-106)]
  \label{def:airline-change-change-needs-choice}
  \uses{def:airline-change-follow-up-action, def:airline-change-involuntary-change}
  \notready
  The airline may change departure times, flight numbers, aircraft,
  operating carrier or itinerary, or cancel a flight, and apply the change
  without the customer. For each change it decides whether to attach
  follow-up actions or send a notice only.

  Case to check: In the worked examples a 10-minute departure change and a
  change of operating carrier carry no follow-up action. A customer moved to
  another carrier has no recorded way to refuse.

  Open point: The sources give no threshold. The business must decide which
  change types and sizes oblige the airline to offer Accept, Cancel or
  Reshop.

  Sources: \href{https://standards.atlassian.net/wiki/spaces/EASD/pages/815202352}{EASD guide, Airline notifies Seller of changes to an Order};
  \href{https://standards.atlassian.net/wiki/spaces/EASD/pages/574455903}{EASD guide, Notification of Airline Initiated Change to an Order}.
\end{definition}

\begin{definition}[Voluntary and involuntary change (C-107)]
  \label{def:airline-change-involuntary-change}
  \uses{def:principal, def:order}
  \notready
  An order change is voluntary when the customer requests it and involuntary
  when the airline imposes it. Resolution 797 commits One Order to support
  servicing, delivery and accounting for both kinds.

  Open point: No source defines the two kinds. Change notices also carry
  changes the customer made directly with the airline, such as a name or
  contact update, so the business must decide which change reasons count as
  involuntary and trigger customer protections.

  Sources: \href{https://www.iata.org/contentassets/72cbd60393ff42b5975d90ce9e049a7d/oneorder-resolution-797.pdf}{IATA Resolution 797};
  \href{https://standards.atlassian.net/wiki/spaces/EASD/pages/672825400}{EASD guide, Inform of Impacted Dependent Services};
  \href{https://retailing.iata.org/armi/docs/ORDOC2/}{IATA capability ORDOC2};
  \href{https://standards.atlassian.net/wiki/spaces/EASD/pages/815202352}{EASD guide, Airline notifies Seller of changes to an Order}.
\end{definition}

\begin{theorem}[Pending follow-up actions (C-108)]
  \label{thm:airline-change-pending-actions}
  \uses{def:airline-change-follow-up-action, def:airline-change-inaction-consequence}
  \notready
  A follow-up action the seller has not performed appears in every order
  view until it is resolved. Performing any one of the actions offered for a
  change resolves all of them.

  Case to check: The itinerary example offers Accept, Cancel, Reshop and
  ContactAirline under one time limit with Auto-acceptance. The seller
  phones the airline, nothing in the order records it, and auto-acceptance
  fires.

  Open point: The source says pending actions are shown but not when they
  stop being pending. The business must decide whether one response clears
  the rest, and whether ContactAirline, which happens offline, ever clears a
  pending Accept.

  Sources: \href{https://standards.atlassian.net/wiki/spaces/EASD/pages/574849094}{EASD guide, Seller Follow-up Actions};
  \href{https://standards.atlassian.net/wiki/spaces/EASD/pages/815202352}{EASD guide, Airline notifies Seller of changes to an Order}.
\end{theorem}

\begin{definition}[Repricing after waitlist confirmation (C-109)]
  \label{def:airline-change-reprice-and-accept}
  \uses{def:airline-change-follow-up-action, def:airline-change-inaction-consequence, def:c08-delivery-status, def:c11-payment-time-limit}
  \notready
  When a waitlisted service is confirmed, the airline may ask the seller to
  reprice the order item and accept. The guide's example pairs this with
  Auto-acceptance if the seller does not act.

  Case to check: The waitlist clears at a higher fare. The seller misses the
  deadline, auto-acceptance confirms the service, and the customer is bound
  to a price never quoted.

  Open point: The source does not say at what price an auto-accepted
  waitlist confirmation stands, or whether the customer must pay a
  difference before the service becomes \texttt{READY TO PROCEED}.

  Source: \href{https://standards.atlassian.net/wiki/spaces/EASD/pages/574849094}{EASD guide, Seller Follow-up Actions}.
\end{definition}

\begin{definition}[Who is told of a change (C-110)]
  \label{def:airline-change-notice-chain}
  \uses{def:principal, def:order}
  \notready
  A service provider reports every service status change to the airline that
  owns the order. That airline notifies the seller of airline-initiated
  changes and the seller acknowledges. The seller is expected to inform the
  customer when a change needs action.

  Case to check: Agency A creates the order and agency B services it later.
  The disruption notice goes only to A, and B keeps selling changes against
  an old version.

  Open point: The sources do not say which seller is notified when several
  sellers have serviced the order, whether the airline must reach the
  customer when the seller does not acknowledge, or whether the seller's
  duty to inform the customer is binding.

  Sources: \href{https://standards.atlassian.net/wiki/spaces/EASD/pages/574455903}{EASD guide, Notification of Airline Initiated Change to an Order};
  \href{https://standards.atlassian.net/wiki/spaces/EASD/pages/672825400}{EASD guide, Inform of Impacted Dependent Services};
  \href{https://standards.atlassian.net/wiki/spaces/EASD/pages/574652557}{EASD guide, Order History};
  \href{https://retailing.iata.org/armi/docs/STPNTF/}{IATA capability STPNTF};
  \href{https://www.iata.org/contentassets/72cbd60393ff42b5975d90ce9e049a7d/oneorder-resolution-797.pdf}{IATA Resolution 797}.
\end{definition}

\begin{theorem}[Cancelled service in an active item (C-111)]
  \label{thm:airline-change-cancelled-in-active-item}
  \uses{def:c06-item-status, def:c08-delivery-status, def:c18-change}
  \notready
  A service with status \texttt{CANCELLED} belongs to an order item with
  status \texttt{CANCELLED}.

  Case to check: In the flight cancellation worked example the second
  segment's flight and bag services become \texttt{CANCELLED} while their
  order item gets no status change.

  Open point: The Order concept ties a \texttt{CANCELLED} service to a
  cancelled item, and the cancellation page treats cancelling some services
  of an item as a change. The flight cancellation example cancels one
  segment's services and leaves the item untouched, so the business must
  decide which model holds.

  Sources: \href{https://standards.atlassian.net/wiki/spaces/EASD/pages/574849062}{EASD guide, Order};
  \href{https://standards.atlassian.net/wiki/spaces/EASD/pages/574783569}{EASD guide, Cancellation of Order Items};
  \href{https://standards.atlassian.net/wiki/spaces/EASD/pages/815202352}{EASD guide, Airline notifies Seller of changes to an Order}.
\end{theorem}

\begin{theorem}[Price of reshopping after airline change (C-112)]
  \label{thm:airline-change-reshop-after-involuntary}
  \uses{def:c18-change, def:airline-change-follow-up-action, def:airline-change-involuntary-change}
  \notready
  A customer who declines an airline-initiated change through the Reshop
  follow-up action pays no change fee.

  Case to check: The new item in the itinerary example carries a 40 EUR
  change fee. A natural reading charges it when the customer reshops away
  from flights the airline moved.

  Open point: The guide says such a reshop may or may not lead to additional
  payment and gives no rule. The business must decide whether change fees
  and fare differences apply, and which alternatives the airline must offer.

  Sources: \href{https://standards.atlassian.net/wiki/spaces/EASD/pages/672825400}{EASD guide, Inform of Impacted Dependent Services};
  \href{https://standards.atlassian.net/wiki/spaces/EASD/pages/574849094}{EASD guide, Seller Follow-up Actions};
  \href{https://standards.atlassian.net/wiki/spaces/EASD/pages/815202352}{EASD guide, Airline notifies Seller of changes to an Order}.
\end{theorem}

\begin{definition}[Seller follow-up action (C-113)]
  \label{def:airline-change-follow-up-action}
  \uses{def:order, def:instant}
  A change notice carries zero or more seller follow-up actions drawn from
  Accept, Reshop, Cancel, ContactAirline and RepriceAndAccept. Each may
  carry a time limit, a consequence of inaction, and an airline contact
  link. A notice with no follow-up action needs only acknowledgement.

  Sources: \href{https://standards.atlassian.net/wiki/spaces/EASD/pages/574849094}{EASD guide, Seller Follow-up Actions};
  \href{https://standards.atlassian.net/wiki/spaces/EASD/pages/574455903}{EASD guide, Notification of Airline Initiated Change to an Order}.
\end{definition}

% Candidate rules for delivery and fulfilment.  Written by hand; see docs/one-order-coverage.md.

\section{Delivery and fulfilment}

\begin{definition}[Delivery status (C-08)]
  \label{def:c08-delivery-status}
  \uses{def:c01-structure, def:delivery-state, def:step}
  \notready
  The delivery status of a service and the moves between statuses.  A service
  is \texttt{CONFIRMED} while the airline waits for payment, becomes
  \texttt{READY TO PROCEED} once paid, and ends \texttt{DELIVERED} or in
  another final status.  This replaces the four delivery states in the
  current model.

  Open point: the source describes thirteen statuses but does not list every
  allowed move or say which statuses are final.

  Sources: \href{https://standards.atlassian.net/wiki/spaces/EASD/pages/574849062/Order}{EASD guide, Order};
  \href{https://retailing.iata.org/armi/docs/ORDSTS/}{IATA capability ORDSTS}.
\end{definition}

\begin{theorem}[Suspended services (C-09)]
  \label{thm:c09-suspended}
  \uses{def:c08-delivery-status}
  A suspended service can only return to the status it had before it was
  suspended, or be removed.

  Source: \href{https://standards.atlassian.net/wiki/spaces/EASD/pages/574849062/Order}{EASD guide, Order}.
\end{theorem}

\begin{definition}[Refunds while delivery is in progress (C-10)]
  \label{def:c10-in-progress}
  \uses{def:c08-delivery-status, def:accept-refund}
  \notready
  A service whose delivery is \texttt{IN PROGRESS} may block refunds and other
  financial actions.

  Open point: the source leaves this to each airline.  The business decides
  whether these rules block them.

  Source: \href{https://standards.atlassian.net/wiki/spaces/EASD/pages/574849062/Order}{EASD guide, Order}.
\end{definition}

% Candidate rules for payment.  Written by hand; see docs/one-order-coverage.md.

\section{Payment}

\begin{definition}[Payment time limit (C-11)]
  \label{def:c11-payment-time-limit}
  \uses{def:c08-delivery-status, def:instant}
  \notready
  A confirmed service waits for payment until the payment time limit the
  airline gave the seller.

  Open point: the source does not say what happens to the service when the
  limit passes without payment.

  Source: \href{https://standards.atlassian.net/wiki/spaces/EASD/pages/574849062/Order}{EASD guide, Order}.
\end{definition}

\begin{definition}[Committing to pay nothing (C-12)]
  \label{def:c12-zero-commitment}
  \uses{def:c08-delivery-status, def:c01-structure}
  \notready
  An order item with nothing to pay goes ahead when the seller commits to pay
  zero.

  Open point: the source does not say which delivery status the item's
  services move to.

  Source: \href{https://standards.atlassian.net/wiki/spaces/EASD/pages/2593488900/Commitment+to+pay+when+no+amount+due}{EASD guide, Commitment to pay when no amount due}.
\end{definition}

\begin{definition}[Paying with several methods (C-13)]
  \label{def:c13-several-methods}
  \uses{def:payment}
  Several payment methods can pay for one order item in a single payment.  The
  payment succeeds only if every method succeeds; if any fails, the whole
  payment fails.

  Source: \href{https://standards.atlassian.net/wiki/spaces/EASD/pages/2592899087/Use+of+multiple+payment+methods+to+pay+for+a+single+item}{EASD guide, Use of multiple payment methods}.
\end{definition}

\begin{theorem}[Amount charged across methods (C-14)]
  \label{thm:c14-amount-covers-price}
  \uses{def:c13-several-methods, thm:c04-item-unsplit}
  \notready
  The amounts charged across the methods cover the item's price plus any
  surcharges.

  Open point: the source's example adds a card surcharge on top of the price
  and states no rule, so it is unclear whether an item paid in full means the
  price was paid or the price plus surcharges.

  Source: \href{https://standards.atlassian.net/wiki/spaces/EASD/pages/2592899087/Use+of+multiple+payment+methods+to+pay+for+a+single+item}{EASD guide, Use of multiple payment methods}.
\end{theorem}

\begin{definition}[Refunds across payment methods (C-17)]
  \label{def:c17-refund-methods}
  \uses{def:c13-several-methods, thm:c14-amount-covers-price, def:accept-refund, thm:refunds-le-payments}
  \notready
  A refund goes back to the payment methods that paid.

  Open point: the source does not say how a refund is split across several
  methods, or whether surcharges are refunded.  If surcharges are refunded, a
  refund can exceed the item's price, which bears on the open question about
  refunds and payments.

  Sources: \href{https://standards.atlassian.net/wiki/spaces/EASD/pages/2592899087/Use+of+multiple+payment+methods+to+pay+for+a+single+item}{EASD guide, Use of multiple payment methods};
  \href{https://standards.atlassian.net/wiki/spaces/EASD/pages/574783569/Cancellation+of+Order+Items}{EASD guide, Cancellation of Order Items}.
\end{definition}

% Second survey round: candidates written from the survey notes in docs/one-order-survey.md.

\begin{definition}[Paid order item (C-126)]
  \label{def:payment-paid-order-item}
  \uses{def:payment, def:c01-structure, def:c08-delivery-status, thm:c14-amount-covers-price, def:c18-change}
  \notready
  An order item is paid once the successful payments allocated to it,
  together with any value carried over from the item it replaced, cover its
  amount due. A \texttt{PENDING} payment does not count. After a change, the
  original payment stays on record and its allocation moves to the replacing
  item.

  Case to check: In the change worked example the new item costs 910.00 and
  the only payment made after the change is 60.00. The item is paid only if
  the 1000.00 payment made for the old item counts toward it.

  Open point: The source speaks of unpaid, partly paid and paid items but
  never states the test. The business must decide whether surcharges count
  toward the amount due and how value carried over from a replaced item is
  counted.

  Sources: \href{https://standards.atlassian.net/wiki/spaces/EASD/pages/574849062/Order}{EASD guide, Order};
  \href{https://retailing.iata.org/armi/docs/PAYORD/}{IATA capability PAYORD};
  \href{https://standards.atlassian.net/wiki/spaces/EASD/pages/842301447/Customer+requests+a+change+of+a+paid+order+prior+to+travel+resulting+in+additional+collection+and+a+refund}{EASD guide, Customer requests a change of a paid order prior to travel resulting in additional collection and a refund};
  \href{https://standards.atlassian.net/wiki/spaces/ARMI/pages/584318980/General+Guidance}{ARMI, General Guidance}.
\end{definition}

\begin{definition}[Payment status and its moves (C-127)]
  \label{def:payment-payment-status}
  \uses{def:payment, def:instant}
  \notready
  A payment transaction is \texttt{PENDING}, \texttt{SUCCESSFUL} or
  \texttt{FAILED}. \texttt{SUCCESSFUL} means the airline is certain to
  collect the funds, \texttt{FAILED} means it knows it will not, and
  \texttt{PENDING} means it cannot yet tell; a redirected payment starts
  \texttt{PENDING} and later moves to \texttt{SUCCESSFUL} or
  \texttt{FAILED}. The moment a payment becomes \texttt{SUCCESSFUL} is its
  payment commitment time.

  Case to check: A settlement plan payment is \texttt{SUCCESSFUL} at order
  time, then the seller remits only 75 percent and the platform marks the
  clearances Not Settled or claws them back. The order still shows the
  payment as \texttt{SUCCESSFUL}.

  Open point: The source does not say whether \texttt{SUCCESSFUL} and
  \texttt{FAILED} are final. The business must decide whether a
  \texttt{SUCCESSFUL} payment can later fail, for example after a card
  chargeback or a seller's default on a settlement plan, and what then
  happens to the items it paid.

  Sources: \href{https://standards.atlassian.net/wiki/spaces/EASD/pages/574849161/Payment+Codes}{EASD guide, Payment Codes};
  \href{https://standards.atlassian.net/wiki/spaces/EASD/pages/574423154/Payment+Status}{EASD guide, Payment Status};
  \href{https://standards.atlassian.net/wiki/spaces/EASD/pages/574455927/Pay+Offline+Using+Redirection}{EASD guide, Pay Offline Using Redirection};
  \href{https://standards.atlassian.net/wiki/spaces/SWO/pages/574653091/Enhanced+and+Simplified+Distribution}{Settlement with Orders, Enhanced and Simplified Distribution};
  \href{https://standards.atlassian.net/wiki/spaces/SWO/pages/574718145/Remittance+of+Funds}{Settlement with Orders, Remittance of Funds}.
\end{definition}

\begin{theorem}[Effects of a payment outcome (C-128)]
  \label{thm:payment-payment-outcome-effects}
  \uses{def:payment-payment-status, def:payment-paid-order-item, def:c08-delivery-status, def:c11-payment-time-limit}
  \notready
  When a payment covering unpaid order items becomes \texttt{SUCCESSFUL},
  every service in those items moves from \texttt{CONFIRMED} to
  \texttt{READY TO PROCEED} and the items lose their payment time limits.
  When it becomes \texttt{FAILED}, the service statuses and payment time
  limits stay as they were.

  Open point: The source states both outcomes for redirected payments only.
  The business must confirm they apply to card and settlement plan payments,
  and decide what a successful payment that covers part of an item's amount
  does to its services.

  Sources: \href{https://standards.atlassian.net/wiki/spaces/EASD/pages/574455927/Pay+Offline+Using+Redirection}{EASD guide, Pay Offline Using Redirection};
  \href{https://standards.atlassian.net/wiki/spaces/EASD/pages/670203914/Payment+Via+Redirection+with+Airline+Notification+of+Payment+Success}{EASD guide, Payment Via Redirection with Airline Notification of Payment Success​};
  \href{https://standards.atlassian.net/wiki/spaces/EASD/pages/574849062/Order}{EASD guide, Order};
  \href{https://standards.atlassian.net/wiki/spaces/EASD/pages/670203914}{EASD guide, Payment Via Redirection with Airline Notification of Payment Success​};
  \href{https://standards.atlassian.net/wiki/spaces/OA/pages/844202130}{Order Accounting, Customer requests a change of a paid order prior to travel resulting in additional collection and a refund}.
\end{theorem}

\begin{theorem}[Time limit on unpaid items (C-129)]
  \label{thm:payment-time-limit-on-unpaid}
  \uses{def:c11-payment-time-limit, def:payment-paid-order-item, def:instant}
  \notready
  Every order item that is unpaid or partly paid carries a payment time
  limit, and a paid item carries none. A time limit of zero on an offer item
  means the item must be paid when the order is created.

  Case to check: An order created without payment returns no time limit for
  one item. Under the second reading that item never expires and holds its
  seat with nothing paid.

  Open point: The ARMI guidance expects the airline to set a limit on every
  unpaid item, while the EASD Time Limits page says a missing limit only
  means no advice was given. The business must decide whether an unpaid item
  may exist without a limit and, if so, whether it can stay unpaid
  indefinitely.

  Sources: \href{https://retailing.iata.org/armi/docs/PAYORD/}{IATA capability PAYORD};
  \href{https://standards.atlassian.net/wiki/spaces/EASD/pages/574423066/Time+Limits}{EASD guide, Time Limits};
  \href{https://standards.atlassian.net/wiki/spaces/EASD/pages/574455879/Create+Order+Without+Payment}{EASD guide, Create Order Without Payment};
  \href{https://standards.atlassian.net/wiki/spaces/EASD/pages/574455927/Pay+Offline+Using+Redirection}{EASD guide, Pay Offline Using Redirection};
  \href{https://retailing.iata.org/armi/docs/ORDWPM/}{IATA capability ORDWPM};
  \href{https://standards.atlassian.net/wiki/spaces/EASD/pages/574849062/Order}{EASD guide, Order}.
\end{theorem}

\begin{definition}[Unpaid item at its time limit (C-130)]
  \label{def:payment-time-limit-expiry}
  \uses{def:c11-payment-time-limit, def:c08-delivery-status, def:c06-item-status, def:payment-payment-status}
  \notready
  When an order item's payment time limit passes without a successful
  payment, the item is removed from the order, and a payment still
  \texttt{PENDING} for it is marked \texttt{FAILED}. The airline may move
  the limit when it accepts a redirected payment as \texttt{PENDING}.

  Case to check: A customer completes a redirected card payment one minute
  after the limit. The airline has already marked the payment
  \texttt{FAILED} and removed the item, so the customer has paid for nothing
  unless the capture is refunded or the item restored.

  Open point: This partly answers the open point of C-11: the Time Limits
  page says removal is expected but names no resulting statuses. The
  business must decide the item status and delivery status after removal
  (for example \texttt{CANCELLED} with \texttt{REMOVED}, or
  \texttt{EXPIRED}) and what happens when a redirected payment is captured
  after the limit.

  Sources: \href{https://standards.atlassian.net/wiki/spaces/EASD/pages/574423066/Time+Limits}{EASD guide, Time Limits};
  \href{https://standards.atlassian.net/wiki/spaces/EASD/pages/574455927/Pay+Offline+Using+Redirection}{EASD guide, Pay Offline Using Redirection};
  \href{https://standards.atlassian.net/wiki/spaces/EASD/pages/1404665857}{EASD guide, 24.1 Release Notes}.
\end{definition}

\begin{theorem}[No payment after the limits (C-131)]
  \label{thm:payment-no-late-payment}
  \uses{def:c11-payment-time-limit, thm:accept-rejects-expired, def:instant}
  \notready
  The airline accepts a payment for an order item only before the item's
  payment time limit and, when the item has a price guarantee, before that
  guarantee expires. A payment attempted after the price guarantee has
  expired fails with a price guarantee error.

  Case to check: The airline extends an item's price guarantee two days past
  its payment time limit, which the source allows. The item can no longer be
  paid while its price is still guaranteed.

  Open point: The payment time limit part is clear, but for the price
  guarantee the source also says repricing is only recommended. The business
  must decide whether an expired guarantee blocks payment or only lets the
  airline change the price, and whether a guarantee may outlast the payment
  time limit, which the source says it should not.

  Sources: \href{https://standards.atlassian.net/wiki/spaces/EASD/pages/574226556/Pay+for+an+Existing+Order}{EASD guide, Pay for an Existing Order};
  \href{https://standards.atlassian.net/wiki/spaces/EASD/pages/574587020/Payment+Errors}{EASD guide, Payment Errors};
  \href{https://standards.atlassian.net/wiki/spaces/EASD/pages/574423066/Time+Limits}{EASD guide, Time Limits}.
\end{theorem}

\begin{theorem}[Failed payment leaves items unpaid (C-132)]
  \label{thm:payment-failed-payment-all-items}
  \uses{def:c13-several-methods, def:payment-paid-order-item}
  One payment transaction may cover several order items with several payment
  methods. If any method's transaction fails, the whole payment fails and
  none of the items it covers becomes paid.

  Case to check: A seller who wants item 1 paid even when the card for item
  2 is declined must send two payments. One combined payment leaves both
  items unpaid.

  Sources: \href{https://standards.atlassian.net/wiki/spaces/EASD/pages/2592899087/Use+of+multiple+payment+methods+to+pay+for+a+single+item}{EASD guide, Use of multiple payment methods to pay for a single item};
  \href{https://standards.atlassian.net/wiki/spaces/EASD/pages/2593325059/Pay+using+multiple+payment+methods+for+a+single+item}{EASD guide, Pay using multiple payment methods for a single item};
  \href{https://standards.atlassian.net/wiki/spaces/EASD/pages/2593030164/Payment+Simulation}{EASD guide, Payment Simulation}.
\end{theorem}

\begin{definition}[Allocating a payment across items (C-133)]
  \label{def:payment-allocation-across-items}
  \uses{thm:payment-failed-payment-all-items, def:payment-paid-order-item, def:c17-refund-methods}
  \notready
  When one payment with several methods covers several order items, the
  airline decides how much of each method goes to each item and the seller
  has no say. A seller who wants a fixed split sends one payment per item.

  Case to check: Items of 800.00 and 100.00 are paid with 400.00 by
  settlement plan and 500.00 by card. If the 100.00 item is cancelled, its
  refund goes to the card under one allocation and to the settlement plan
  under another.

  Open point: The source gives the airline the allocation and states no rule
  for it. The business must fix the rule, because it decides which method is
  refunded when one of the items is later cancelled.

  Sources: \href{https://standards.atlassian.net/wiki/spaces/EASD/pages/2593030164/Payment+Simulation}{EASD guide, Payment Simulation};
  \href{https://standards.atlassian.net/wiki/spaces/EASD/pages/2592899087/Use+of+multiple+payment+methods+to+pay+for+a+single+item}{EASD guide, Use of multiple payment methods to pay for a single item};
  \href{https://retailing.iata.org/armi/docs/PAYMIX/}{IATA capability PAYMIX}.
\end{definition}

\begin{definition}[Combinable payment methods (C-134)]
  \label{def:payment-combinable-methods}
  \uses{def:c13-several-methods}
  \notready
  The payment catalogue places each accepted method in one or more groups,
  and two methods may be combined in one payment when they share a group.
  The airline refuses a payment that combines methods it does not accept
  together.

  Case to check: Method A is in groups 1 and 2, B in groups 2 and 3, and C
  in groups 1 and 3. Every pair shares a group, but no group holds all
  three.

  Open point: The source defines combinability for pairs only. The business
  must decide whether three or more methods need one group that holds them
  all or only a shared group for each pair.

  Sources: \href{https://standards.atlassian.net/wiki/spaces/EASD/pages/2592899087/Use+of+multiple+payment+methods+to+pay+for+a+single+item}{EASD guide, Use of multiple payment methods to pay for a single item};
  \href{https://standards.atlassian.net/wiki/spaces/EASD/pages/2593325059/Pay+using+multiple+payment+methods+for+a+single+item}{EASD guide, Pay using multiple payment methods for a single item}.
\end{definition}

\begin{definition}[Surcharge on each payment method (C-135)]
  \label{def:payment-surcharge-per-method}
  \uses{thm:c14-amount-covers-price, def:c13-several-methods, def:money-add}
  Each method in a payment has a requested amount and a surcharge that
  depends on the method and that amount, and the amount charged to the
  method is the requested amount plus its surcharge. The requested amounts
  of one payment add up to the whole cost of the items it pays.

  Case to check: Reading paid in full as charged amounts equal to the price
  fails on the source's own example: a 100.00 item is paid with 20.00 by
  settlement plan and 82.00 by card, of which 2.00 is the card surcharge.
  This settles most of the open point of C-14.

  Sources: \href{https://standards.atlassian.net/wiki/spaces/EASD/pages/2593030164/Payment+Simulation}{EASD guide, Payment Simulation};
  \href{https://standards.atlassian.net/wiki/spaces/EASD/pages/2592899087/Use+of+multiple+payment+methods+to+pay+for+a+single+item}{EASD guide, Use of multiple payment methods to pay for a single item}.
\end{definition}

\begin{theorem}[Quoted surcharge at payment time (C-136)]
  \label{thm:payment-quoted-surcharge}
  \uses{def:payment-surcharge-per-method, thm:c14-amount-covers-price}
  \notready
  When the airline has quoted a precise surcharge for a stated card brand
  and card type, the payment must use that brand and type and is charged
  that surcharge. A payment with a different brand, scheme, or consumer or
  commercial card type fails with a mismatch error.

  Case to check: The offer is priced with a 3.00 surcharge for a consumer
  Visa debit card, and at payment the airline charges 5.00 for the same
  card. Nothing in the source forbids it.

  Open point: The mismatch errors exist, but no page says a quoted surcharge
  binds the airline, and the catalogue calls its ranges estimates. The
  business must decide whether the charged surcharge may differ from the
  quote or fall outside a published minimum and maximum.

  Sources: \href{https://standards.atlassian.net/wiki/spaces/EASD/pages/574488711/Payment+Catalogue}{EASD guide, Payment Catalogue};
  \href{https://standards.atlassian.net/wiki/spaces/EASD/pages/574685328/Payment+Preference}{EASD guide, Payment Preference};
  \href{https://standards.atlassian.net/wiki/spaces/EASD/pages/574587020/Payment+Errors}{EASD guide, Payment Errors};
  \href{https://standards.atlassian.net/wiki/spaces/ARMI/pages/584318980/General+Guidance}{ARMI, General Guidance}.
\end{theorem}

\begin{definition}[Refundable payment surcharges (C-137)]
  \label{def:payment-surcharge-refundability}
  \uses{def:payment-surcharge-per-method, def:c17-refund-methods, def:accept-refund, thm:refunds-le-payments}
  \notready
  Each payment surcharge carries an indicator saying whether it is
  refundable. A refund of an item's payment returns a surcharge only when
  its indicator marks it refundable.

  Case to check: A 100.00 item is paid with 20.00 by settlement plan and
  82.00 by card, including a non-refundable 2.00 surcharge. A full refund of
  102.00 returns the surcharge, while the cap under the indicator is 100.00.

  Open point: The indicator appears in every simulation example, always
  marked non-refundable, and no page says how a refund uses it. The business
  must decide whether the indicator binds refunds; with refundable
  surcharges a refund can exceed the item's price, which bears on C-17 and
  the open refund question.

  Sources: \href{https://standards.atlassian.net/wiki/spaces/EASD/pages/2593030164/Payment+Simulation}{EASD guide, Payment Simulation};
  \href{https://standards.atlassian.net/wiki/spaces/EASD/pages/2592899087/Use+of+multiple+payment+methods+to+pay+for+a+single+item}{EASD guide, Use of multiple payment methods to pay for a single item}.
\end{definition}

\begin{definition}[Refund as a negative payment (C-138)]
  \label{def:payment-refund-negative-payment}
  \uses{def:refund, def:payment, def:totals, thm:refunds-le-payments, def:c17-refund-methods}
  \notready
  A refund appears in the order as a payment transaction of its own, with a
  new payment ID, a negative amount, a payment method and its own status. In
  the worked examples the method is the one that paid, but the record does
  not point to the payment it returns.

  Case to check: An order holds a 1000.00 settlement plan payment.
  Cancelling it adds a payment of minus 1000.00 with its own ID and no
  reference to the first.

  Open point: The source does not link a refund to the payment it reverses,
  while our model attaches each refund to one payment. The business must
  decide whether to keep that link; without it the per payment limit in the
  open refund question cannot be checked and only an order level limit
  remains.

  Sources: \href{https://retailing.iata.org/armi/docs/PAYREF/}{IATA capability PAYREF};
  \href{https://retailing.iata.org/armi/docs/PAYSUM/}{IATA capability PAYSUM};
  \href{https://standards.atlassian.net/wiki/spaces/EASD/pages/817594417/Customer+requests+a+cancellation+to+paid+Order+or+OrderItems+resulting+in+full+refund}{EASD guide, Customer requests a cancellation to paid Order or OrderItems resulting in full refund};
  \href{https://standards.atlassian.net/wiki/spaces/EASD/pages/842301447/Customer+requests+a+change+of+a+paid+order+prior+to+travel+resulting+in+additional+collection+and+a+refund}{EASD guide, Customer requests a change of a paid order prior to travel resulting in additional collection and a refund}.
\end{definition}

\begin{definition}[Retrying a failed payment (C-139)]
  \label{def:payment-failed-payment-recovery}
  \uses{def:payment-payment-status, thm:payment-payment-outcome-effects, def:c11-payment-time-limit}
  \notready
  After a payment fails, its items stay unpaid with services and time limits
  unchanged, and the seller may try again in a new payment transaction
  before the time limit. The error tells the seller whether to correct the
  card details, add payer authentication, use another method, or stop; after
  a hard stop the airline accepts no payment from that payer.

  Case to check: A payer flagged for suspected fraud on one order pays for a
  second order with the same card. The source does not say whether the stop
  carries over.

  Open point: The source calls recovery optional and does not say whether
  the airline must accept a retry. The business must decide how many retries
  are allowed, whether a hard stop applies to the payer, the order or the
  method, and whether the order is then cancelled.

  Sources: \href{https://standards.atlassian.net/wiki/spaces/EASD/pages/574423154/Payment+Status}{EASD guide, Payment Status};
  \href{https://standards.atlassian.net/wiki/spaces/EASD/pages/574587020/Payment+Errors}{EASD guide, Payment Errors};
  \href{https://standards.atlassian.net/wiki/spaces/EASD/pages/574455927/Pay+Offline+Using+Redirection}{EASD guide, Pay Offline Using Redirection};
  \href{https://standards.atlassian.net/wiki/spaces/EASD/pages/574226627/Payer+Authentication}{EASD guide, Payer Authentication};
  \href{https://retailing.iata.org/armi/docs/PAYRCV/}{IATA capability PAYRCV}.
\end{definition}

\begin{theorem}[Commitment only when nothing due (C-140)]
  \label{thm:payment-commitment-zero-due}
  \uses{def:c12-zero-commitment, def:payment-paid-order-item}
  \notready
  The commitment payment method always carries an amount of zero and may
  cover only items with nothing to pay or collect. Items with nothing due
  may also be covered by a normal payment for a payable item, and every item
  a payment covers, zero priced ones included, is listed on it.

  Open point: The EASD guide added a commitment method in release 26.1,
  while the ARMI guidance and the change worked examples record the same
  commitment as an OT payment of zero. The business must pick one form and
  decide whether a zero OT payment counts when testing whether an item is
  paid.

  Sources: \href{https://standards.atlassian.net/wiki/spaces/EASD/pages/2593488900/Commitment+to+pay+when+no+amount+due}{EASD guide, Commitment to pay when no amount due};
  \href{https://standards.atlassian.net/wiki/spaces/EASD/pages/2533031964/26.1+Release+Notes}{EASD guide, 26.1 Release Notes};
  \href{https://standards.atlassian.net/wiki/spaces/ARMI/pages/584318980/General+Guidance}{ARMI, General Guidance};
  \href{https://standards.atlassian.net/wiki/spaces/EASD/pages/842268768/Customer+requests+a+change+of+a+paid+order+prior+to+travel+resulting+in+a+refund}{EASD guide, Customer requests a change of a paid order prior to travel resulting in a refund}.
\end{theorem}

\begin{definition}[Paying with a voucher (C-141)]
  \label{def:payment-voucher-payment}
  \uses{thm:c20-value-store, def:c17-refund-methods, def:c13-several-methods}
  \notready
  A voucher is a payment method that spends value the airline holds for the
  payer, alone or combined with other methods. Value stored as an order item
  can be spent in another order only after the airline issues it as a
  voucher.

  Case to check: An 800.00 item is paid with a 300.00 voucher and 500.00 by
  card, then cancelled with a 100.00 penalty. Refunding the card first
  returns 500.00 to the card and 200.00 to the voucher; a pro rata split
  returns 437.50 and 262.50.

  Open point: The voucher capability has no guidance beyond its name. The
  business must decide whether a voucher can be partly spent, whether it
  expires, and whether a refund of a voucher paid amount goes back to the
  voucher or to another method.

  Sources: \href{https://retailing.iata.org/armi/docs/PAYVCH/}{IATA capability PAYVCH};
  \href{https://retailing.iata.org/armi/docs/PAYMIX/}{IATA capability PAYMIX};
  \href{https://standards.atlassian.net/wiki/spaces/EASD/pages/574554245/Value+Store}{EASD guide, Value Store};
  \href{https://standards.atlassian.net/wiki/spaces/EASD/pages/574619671/Payment+Concepts}{EASD guide, Payment Concepts}.
\end{definition}

% Candidate rules for commission, accounting and settlement.  Written by hand; see docs/one-order-coverage.md.

\section{Commission, accounting and settlement}

\begin{theorem}[Billing after delivery (C-22)]
  \label{thm:c22-bill-after-delivery}
  \uses{def:c08-delivery-status}
  A service provider does not bill the airline for a service before it has
  been delivered or consumed, unless the two have agreed otherwise.

  Source: \href{https://www.iata.org/contentassets/72cbd60393ff42b5975d90ce9e049a7d/oneorder-resolution-797.pdf}{IATA Resolution 797}.
\end{theorem}

